\documentclass[11pt,letterpaper]{article}
\usepackage[T1]{fontenc}
\usepackage[utf8]{inputenc}
\usepackage{lmodern,microtype}
\ifdefined\pdfmapfile
\pdfmapfile{+lm.map}
\pdfmapfile{+cm.map}
\pdfmapfile{+symbols.map}
\fi
\usepackage[margin=0.9in,headheight=15pt]{geometry}
\usepackage{amsmath,amssymb,amsthm,mathtools}
\usepackage{booktabs,array,longtable,enumitem}
\usepackage{xcolor,listings,fancyhdr,xurl,needspace}
\definecolor{inkblue}{RGB}{26,56,89}
\definecolor{muted}{RGB}{74,84,95}
\usepackage[hypertexnames=false,colorlinks=true,linkcolor=inkblue,citecolor=inkblue,urlcolor=inkblue]{hyperref}
\hypersetup{pdftitle={iPRA with inferred elimination premises},pdfauthor={OpenAI, prepared for Vladimir Reshetnikov},pdfsubject={Complete iPRA proof checkers, equivalence proofs and reproducible complexity comparisons}}
\newtheorem{theorem}{Theorem}[section]
\newtheorem{lemma}[theorem]{Lemma}
\newtheorem{proposition}[theorem]{Proposition}
\newtheorem{corollary}[theorem]{Corollary}
\theoremstyle{definition}
\newtheorem{definition}[theorem]{Definition}
\newtheorem{remark}[theorem]{Remark}
\newcommand{\D}{\mathsf D}
\newcommand{\Q}{\mathsf Q}
\newcommand{\A}{\mathcal A}
\newcommand{\G}{\mathcal G}
\newcommand{\U}{\mathcal U}
\newcommand{\ipra}{\mathrm{iPRA}}
\setlength{\parindent}{0pt}
\setlength{\parskip}{5pt plus 1pt}
\setlength{\emergencystretch}{2em}
\setlist{topsep=4pt,itemsep=2pt,parsep=1pt}
\lstset{basicstyle=\ttfamily\scriptsize,breaklines=true,breakatwhitespace=false,
 columns=fullflexible,keepspaces=true,showstringspaces=false,frame=single,
 rulecolor=\color{muted!30},backgroundcolor=\color{muted!3},
 xleftmargin=3pt,xrightmargin=3pt,aboveskip=7pt,belowskip=7pt}
\pagestyle{fancy}
\fancyhf{}
\fancyhead[L]{\small\color{muted}iPRA with inferred elimination premises}
\fancyhead[R]{\small\color{muted}12 November 2026}
\fancyfoot[C]{\thepage}
\renewcommand{\headrulewidth}{0.3pt}
\title{\color{inkblue}\textbf{iPRA with inferred elimination premises}\\[5pt]
\large Smaller proof checkers below the audited comparison envelope}
\author{Prepared for Vladimir Reshetnikov}
\date{12 November 2026}
\begin{document}
\maketitle
\thispagestyle{empty}
\begin{abstract}
We simplify a complete proof-certificate presentation of intuitionistic
primitive recursive arithmetic by reconstructing an elimination's major
formula from its target and substitution data. The major must still be
proved. The resulting programs reduce the identified compact comparison
envelope from 256 significant Python tokens, 278 AST nodes and 506 normalized
bytes to 253 tokens, 275 nodes and 503 bytes, with different programs attaining
the three coordinates. A separate complete native public-goal checker falls
from 656 to 643 bytes, also improving both other coordinates. We give universal
simulation arguments, executable certificate translations, frozen controls,
and the full inherited logical and arithmetic bi-interpretation. Both
numerical domains are whole, equality is actual, and both composites have
graph $x=y$. The mathematical axioms are inherited; the new results concern
certificate structure and trusted implementation size. No absolute minimum
or new axiom deletion is claimed.
\end{abstract}
\vfill
\begin{center}\small
Comparison snapshot: \texttt{62396b06020f794fa6dc01b4025a6a5db9d4d28a}\\
Standalone LaTeX source, compiled PDF, proof checkers, controls and validation artifacts.
\end{center}
\clearpage
\tableofcontents
\clearpage
\section{Result and its exact scope}

This report gives a smaller implementation of a complete presentation of
intuitionistic primitive recursive arithmetic. The mathematical theory is
the native scalar-iteration calculus already developed in ProveIt. The new
contribution is a simplification of its proof-certificate protocol, together
with two independently justified implementation refinements. No arithmetic
axiom deletion and no theorem of absolute minimality is claimed.

The comparison is pinned to commit
\texttt{62396b06020f794fa6dc01b4025a6a5db9d4d28a}, whose ConciseIPRA
directory contains ninety research packages. Smaller implementation
supplements inside those packages matter: the README table alone does not
give the attained comparison frontier.\cite{upstream} In particular, the byte target is
506, not 521 or 516; the complete native public-goal checker target is 656.
The relevant source files, their hashes, and the measurements are included
in this archive.

\subsection{What has become smaller}

An elimination certificate formerly supplied both a template with its
endpoints and a redundant copy of the major formula. Once the proposed
target instance matches the current goal, that major can be reconstructed.
The verifier still demands a proof of the reconstructed formula. Removing
the comparison against the redundant copy reduces the token and abstract
syntax tree coordinates. This is a deterministic change in the certificate
protocol, rather than an invitation to search for an omitted proof.

A separate implementation represents successful recursive checks by
nonempty result objects and sequences the assumption construction and its
continuation inside an ordered pair. This reduces the normalized byte
coordinate without changing the certificate language. Finally, a small
Boolean-order identity reduces the complete native public-goal validator
while preserving its entire finite-JSON acceptance relation.

\begin{quote}
\textbf{Comparison boundary.} The compact profiles have the inherited
guarded public-language contract and do not separately validate that
public grammar. The native profile performs the inherited full Eq/$\D$
public-goal grammar and scope validation. These two boundaries are measured
separately throughout.
\end{quote}

The three compact coordinate improvements are attained by different
programs. There is no claim that one program simultaneously has the
smallest token, node, and byte counts. Each proposed profile nevertheless
implements the same public theoremhood relation, through its specified
certificate protocol. The native profile gives a separate, stronger input
boundary and improves all three of its incumbent's coordinates in one
program. The full measurements appear in Section~\ref{sec:measurements}.

\subsection{The statement established}

Let $\mathcal N$ be the native calculus specified below, and let
$\mathcal A$ be its root-selected auxiliary calculus restricted to guarded
public conclusions, as specified in Appendix~\ref{sec:composition}.
For the appropriate public encoding of a closed formula $A$:
\begin{enumerate}[label=(\roman*)]
\item the new native checker has an accepting finite certificate exactly
      when the old native checker does;
\item the inferred-major compact checkers have an accepting certificate
      exactly when $\mathcal A$ proves $A$;
\item the byte-profile transformation preserves acceptance of every finite
      certificate for its control, under the same ideal resource convention;
\item the ordinary realizations of these public presentations are
      bi-interpretable and equiconsistent with full first-order iPRA.
\end{enumerate}
Both numerical interpretation domains in (iv) are the whole original
domain, equality is actual equality, and both numerical isomorphisms have
the graph $x=y$. The last statement is inherited through exact proof
translations; this report supplies the new checker equivalences and
reproduces the logical and arithmetic bridge in its appendices. No new
Lean or Rocq formalization is asserted.

\section{The complete native theory}
\label{sec:main-theory}

There is one numerical sort. Terms and formulas are
\[
 t ::= x\mid 0\mid S(t)\mid I(n,a;u.h),\qquad
 A ::= (s=t)\mid\Q_y A\mid\D_{x,y}A.
\]
In an iterator only $u$ in $h$ is bound; its counter and initial value are
ordinary terms. There are no function-valued object variables and no
universal evaluator. A particular iterator is a finite syntactic
description with fixed numerical parameters. Put $1=S0$.

The intended ordinary readings are
\[
 \Q_yB(y)\;\equiv\;B(0)\to\forall y B(y),\qquad
 \D_{x,y}B(x,y)\;\equiv\;
 \exists x\bigl(B(x,0)\to\forall y B(x,y)\bigr).
\]
The designated contradiction in native syntax is $0=1$.
Contexts $\Gamma$ are finite lists of reusable formulas. All substitutions
are capture-free; bound variables are renamed before substitution.

\subsection{Nonlogical axiom schemas}

Exactly two computation schemas are primitive:
\[
 I(0,a;u.h)=a,\qquad
 I(Sn,a;u.h)=h[I(n,a;u.h)/u].
\tag{I0,IS}
\]
They are available in every retained context. In particular, successor
injectivity, zero separation, pairing, predecessor, and an equality test
are not additional primitive numerical axioms of this native presentation.

\subsection{Logical and structural inference rules}

Besides computation, the primitive patterns are the following.
\[
 \frac{A\in\Gamma}{\Gamma\vdash A}\quad\mathrm{Hyp},\qquad
 \frac{\Gamma\vdash s=t\quad\Gamma\vdash F(s)}
      {\Gamma\vdash F(t)}\quad\mathrm{Tr},
\]
\[
 \frac{\Gamma\vdash\Q_yB\quad\Gamma\vdash B(0)}
      {\Gamma\vdash B(t)}\quad\mathrm{QE},
\]
\[
 \frac{\Gamma\vdash\D_{x,y}B\quad
       \Gamma,\Q_yB(b,y)\vdash A}
      {\Gamma\vdash A}\quad\mathrm{DE},
\]
\[
 \frac{\Gamma,B(t,0)\vdash B(t,Sz)}
      {\Gamma\vdash\D_{x,y}B(x,y)}\quad\mathrm{DI}_{\rm case},
\]
\[
 \frac{\Gamma,E(t,z)\vdash E(t,Sz)}
      {\Gamma\vdash\D_{x,y}E(x,y)}\quad\mathrm{DI}_{\rm progress}.
\]
Here $F$ is any native formula template, while $E$ in the last rule must
be a literal equation. In DE the parameter $b$ is fresh for the retained
context, the conclusion, and the fixed parameters of $B$. In either DI
rule, $z$ is fresh for the context, witness $t$, and fixed parameters of the
matrix. The displayed local assumptions are discharged. The two DI
patterns form one introduction family with a checked mode bit.

There are no other primitive logical axioms or inference rules. Weakening,
exchange, contraction, renaming, cut, and numerical substitution are
admissible for this unrestricted native presentation by induction on
finite derivations. They are metatheorems, not uncharged executable
instructions. In presentations with a syntactic source restriction,
substitution must instead be recovered by a separate admissibility
argument; the appendices give that argument where it is needed.

The auxiliary $\Q$ can be eliminated from public formulas: $\Q_yB(y)$ is
interderivable with $\D_{v,y}B(y)$ for an unused $v$. The checkers therefore
use Eq/$\D$ as their native public language while retaining $\Q$ as an
internal resource.

\section{Why this is full intuitionistic arithmetic}
\label{sec:main-bridge}

It is necessary to recover all first-order intuitionistic connectives,
not just an equational or negative fragment. Here are the constructions
and the dependency order; the detailed inherited proofs are reproduced
after the new results.

\subsection{Equation induction and numerical cases}

For a literal equation $E(y)$, introduce $\D_{x,y}E(y)$ by progress with
unused $x$. Eliminate it to obtain the closure $\Q_yE(y)$. Supplying
$E(0)$ and applying QE yields $E(t)$. Thus contextual equation induction is
derived. Cases for an arbitrary formula are obtained in the same way
using case DI, ignoring its local zero assumption when proving the
successor case. Freshness is retained in both arguments.

Reflexivity follows from computation and transport. With
$r=I(0,t;u.0)$, use $r=t$ as both transport premises and the template
$v=t$ to obtain $t=t$. Standard templates then give symmetry,
transitivity, and congruence. Transport under a binder is part of the
native rule; reduced variants recover it explicitly.

\subsection{All six ordinary logical operations}

Let $C(v,a,b)=I(v,a;u.b)$, with unused $u$. Its two computations select
$a$ at zero and $b$ at a successor. Define a gate by
\[
 G_v(s=t)=(C(v,0,s)=C(v,0,t)),\qquad
 G_v(\D_{x,y}B)=\D_{x,y}G_vB.
\]
Then $G_0A$ is derivable and $G_{Sz}A$ is interderivable with $A$.
One also constructs, by external recursion on the finite formula
syntax, a selector $M_v(A,B)$ with endpoints
$M_0(A,B)\dashv\vdash A$ and $M_{Sz}(A,B)\dashv\vdash B$.
For equations it applies $C$ to both sides; when binder depths differ,
the shorter formula is padded by a vacuous gated binder. The complete
recursion is given in Appendix~\ref{subsec:choice-logic}.

Let $p(z)=I(z,0;u.C(u,1,0))$. The formulas
\[
\begin{aligned}
 A\Rightarrow_*B&=\D_{u,z}M_z(A,B),\\
 A\wedge_*B&=\D_{u,z}G_zM_{p(z)}(A,B),\\
 \exists_*x\,A(x)&=\D_{x,y}G_yA(x),\\
 A\vee_*B&=\exists_*x\,M_x(A,B),\\
 \forall_*x\,A(x)&=A(0)\wedge_*\D_{u,z}G_zA(z),\\
 \bot_*&=(0=1)
\end{aligned}
\]
have all the ordinary intuitionistic introduction and elimination
rules. Every newly introduced binder is fresh, and $u$ is unused.
For conjunction, the instances $p(1)=1$ and $p(2)=0$ recover its two
components. For the universal quantifier, the first conjunct supplies
zero and the second supplies every successor; numerical coverage then
supplies every element. No excluded middle for quantified formulas is
used.

\subsection{The arithmetic bridge and both composites}

Define the ordinary theory $T_I$ using intuitionistic first-order equality
logic, the same staged iterator equations, contextual equation induction,
and the two sentences
\[
 \forall x\,(x=0\vee\exists y\,x=Sy),\qquad (0=S0)\to\bot.
\]
The second sentence is essential: a merely positive equational theory
would have a singleton model. In native syntax its translation is the
identity implication $(0=1)\Rightarrow_*(0=1)$, because native falsity is
that equation. This explicitly reconciles ordinary falsity and native
contradiction.

The internal arithmetic proof first constructs a scalar predecessor, then
a characteristic of actual equality, then recovers quantifier-free
induction, and finally constructs pairing and compiles primitive recursion.
This is the dependency order of the proof. In particular, predecessor
is not constructed using an assumed faithful equality detector or an
assumed pairing function. Each fixed primitive-recursive description is
compiled to a scalar iterator using derived pairing to store its counter
and current result. All passive parameters remain fixed during induction.

In the reverse interpretation, iPRA defines every fixed scalar iterator
by primitive recursion that ignores the recursion-counter argument in its
update. Equation induction proves uniqueness of each fixed iterator and
of each primitive recursion. External induction on descriptions therefore
shows, internally in both theories, that a twice-translated operation
equals the original operation at every argument. The numerical domain is
unchanged, so the composite isomorphism is exactly $x=y$.

For formulas the essential composite is
\[
 \D_{x,y}B\dashv\vdash
 \exists_*x\bigl(B(x,0)\Rightarrow_*\forall_*y\,B(x,y)\bigr).
\]
One direction opens the native closure and introduces the three macros;
the other eliminates the existential and introduces $\D$ using the
implication and universal premise. External induction on formulas and
the endpoint properties of $G,M$ give both formula composites. The
explicit falsity bridge then yields effective contradiction-preserving
proof translations and equiconsistency. These arguments do not assert
that either theory proves its own consistency.

\section{The compact public presentation}
\label{sec:public-boundary}

The compact profiles use the same numerical language and a guarded public
formula language $\mathcal G$. Every equation is public. A formula
$\D_{x,y}B$ is public if $B$ is compound and public, or if it is literally
\[
 I(y,0;u.h)=I(y,0;u.k).
\]
The update bodies may contain their state variable and all displayed
parameters. Its zero instance is provable by computation. Auxiliary
proofs may additionally use unrestricted Eq/$\Q$/$\D$ formulas.

In the auxiliary calculus, transport has a source other than the literal
$0$, and one introduction chooses progress for a literal equation and
cases for a compound matrix. This is a complete presentation on
\emph{public} sequents. It is not asserted to be equivalent on every
auxiliary raw formula. The public guard is mathematically relevant, not
an optional output formatting convention.

For clarity, the translation of an arbitrary native formula uses the
padding
\[
 L(E)=\D_{u,v}G_vE,\qquad
 P(A)=\begin{cases}L(A),&A\text{ an equation},\\A,&A\text{ compound}.
 \end{cases}
\]
Padding is interderivable with the original formula. First replace every
auxiliary $\Q_yB$ by $\D_{v,y}B$ with unused fresh $v$. On the resulting
Eq/$\D$ syntax define $T(E)=E$ for equations and
$T(\D_{x,y}B)=\D_{x,y}P(T(B))$. By induction $T(B)$ is public;
its padding is compound and public, so the outer binder is also public.
The guarded arithmetic basis recovers equation induction, numerical
cases, and both directions of padding. The inherited source-restriction
argument restores zero-source transport without adding an axiom. The
appendices prove the translation and its formula round trips.

\subsection{Finite array representation}

The certificate is a JSON array. Its last element is the goal; preceding
four-field packets are consumed from the end. An unused prefix is allowed
by the inherited contract. Success is zero process exit status together
with exactly \texttt{ACCEPT} on standard output.

In the uniform arity codec a unary array represents a paired binder,
a binary array an equation, and a ternary array an internal closure:
\[
 [B]\leftrightarrow\D_{x,y}B,\qquad
 [s,t]\leftrightarrow s=t,\qquad
 [B,B,B]\leftrightarrow\Q_yB.
\]
In numerical position, a unary array is successor and a ternary array is
the initialized scalar iterator. Array nodes adjust the opening depth by
$5-\operatorname{len}(a)$; paired numerical slots have width two.
Negative integer labels are allocated eigenparameters. The scope floor
decreases before the continuation opens its fresh assumption.

The representation contains administrative truth and zero aliases.
Their complete sorted reading, the dummy slots in successor and in the
iterator's counter and seed, and the known-zero-instance supply lemma
are reproduced in Appendix~\ref{sec:arity-checker}. No lemma equating
arbitrary raw spellings is assumed. Public meanings and proof-level
administrative objects remain distinct in the correctness statement.

\subsection{The current seed-selected computation packets}

The inherited arity appendix records an older packet arrangement. The
measured controls use the following equivalent seed-selected arrangement,
also documented verbatim in \texttt{reference/packet\_FORMAT.md}. A packet
is $[m,s,u,F]$. After any transfer, a nonzero major is proved and its
triplication is retained: a unary paired formula supplies its ternary
closure, while the other triplications are administrative truth entries.
If the major is falsy and the initializer field $u$ is truthy, the generated
computation equation is
\[
 \begin{cases}
 I(0,u;v.F)=u,&s\text{ falsy},\\
 I(Ss,u;v.F)=F[I(s,u;v.F)/v],&s\text{ truthy}.
 \end{cases}
\]
If both $m$ and $u$ are falsy, the current unary goal is introduced with
witness $F$. The native profile's falsy field distinguishes the two
permitted introduction modes; the compact profile reads the mode from
the matrix root. The common continuation then allocates a fresh parameter
and checks the successor goal under the selected opened assumption.

This arithmetic dispatch does not remove the computation equations whose
initializer or successor predecessor is literally zero. Put
$W=I(2,1;v.0)$. Its successor computation is an allowed direct packet and
proves $W=0$. Given an omitted computation instance, replace just its
forbidden zero initializer and/or predecessor by the same placeholder
$w$, leaving passive zeros in its update unchanged. The resulting
equation $E(W)$ is an allowed direct computation. Transport along $W=0$
proves $E(0)$, the required instance. Thus two computation leaves and one
transport recover every omitted instance, without deciding any term's
numerical value.

Generated assumptions can contain the newly admitted parameter and need
not be uniformly true in it. The precise justification is the
known-zero-instance supply lemma: if $b$ is fresh for $\Gamma,A$,
$H(0)$ is proved, and $\Gamma,H(b)\vdash A$, substitute zero for $b$
in the continuation and cut in $H(0)$. The sorted opener proof establishes
that a successful raw computation supply has such a zero instance.
This is why changing the order of scope allocation and opening would
require a new proof. The transformations below preserve that order.

\section{New result: infer the elimination major}
\label{sec:inferred}

Write $i(F,t)$ for the existing capture-safe opening of the raw template
$F$ at the numerical term $t$, at the current scope floor. A packet has
fields $(m,s,u,F)$. The old transfer test was
\[
 m=M(s,u,F)\quad\hbox{and}\quad i(F,u)=A,
\]
where
\[
 M(s,u,F)=
 \begin{cases}
 [s,u],&s\text{ is truthy in the codec},\\
 [F,F,F],&s\text{ is falsy in the codec}.
 \end{cases}
\]
The discriminator is certificate truthiness (in particular, an empty
array is falsy), not a decision about the arithmetic value of a term.
On canonical payloads it distinguishes the literal zero from other term
constructors. The first alternative requires equality transport and the
second closure specialization. In either case the verifier checks the major and then
continues with $i(F,s)$.

The new protocol makes the target match decisive:
\[
 \boxed{\ i(F,u)=A\quad\Longrightarrow\quad
       m:=M(s,u,F),\quad A:=i(F,s).\ }
\]
The reconstructed $m$ is subsequently proved by exactly the same recursive
major-premise checker. A supplied packet field cannot assert the
reconstructed premise without proving it. When the target does not match,
the packet retains the old interpretation.

\begin{theorem}[Soundness by effective-major expansion]
Every accepting inferred-major certificate over the documented finite
integer/list language expands to an accepting certificate for the matched
old checker, proving the same public goal.
\end{theorem}
\begin{proof}
Follow the finite successful execution. At each packet at which the target
matches, replace its stored major by $M(s,u,F)$; at other packets leave it
unchanged. Copy the successful major subexecution and the continuation
recursively. The old checker now passes both comparisons, constructs the
same assumption, and performs the same openings. If the target fails, its
old conjunction fails and neither protocol transfers. The scope floor,
context, and packet-consumption order therefore agree by induction on the
finite execution tree. The transformation is implemented outside the
trusted checker in \texttt{translate\_inferred.py}.
\end{proof}

\begin{lemma}[Separation of numerical payloads and meaningful goals]
On canonical numerical payloads, $i(F,u)$ evaluated at formula depth zero
cannot equal a meaningful Eq/$\Q$/$\D$ goal.
\end{lemma}
\begin{proof}
A numerical tree has only unary and ternary array nodes. Starting at even
formula depth, their depth increments are respectively four and two, so
depth remains even. At leaves, the opening either preserves a protected
integer or supplies the formula-neutral zero; it does not insert the
replacement image at a meaningful numerical hole. Hence the result still
has no binary node. Every meaningful native formula has a binary equation
node at the bottom of its finite chain of binders. Structural equality
between the two trees is therefore impossible.
\end{proof}

\begin{theorem}[Completeness of inferred majors]
Every theorem of the guarded public calculus has a finite inferred-major
certificate. Thus the new protocol presents the same public theory as its
matched compact control.
\end{theorem}
\begin{proof}
Use the inherited canonical proof compiler. In a transport or closure-use
packet the target matches and the inferred major is exactly its required
major. In a computation or introduction packet the fourth field is a
numerical payload, so the preceding lemma prevents an accidental match.
An ordinary paired elimination chooses the inert payload $F=0$; its goal
is meaningful and cannot equal that scalar. Hypothesis closure consumes no
packet. These cases cover the compiled derivation. In the local profile
neutral majors close using the retained scalar zero; the global profile
skips proof recursion for a falsy major. The existence of the old canonical
compiler and the guarded interpretation
then establish completeness, rather than an empirical assertion about
the regression suite.
\end{proof}

This theorem is deliberately not a claim of identical acceptance on every
old raw certificate. The protocol now accepts some packets with a
redundant major field different from the inferred major. The expansion
proof establishes soundness, and the canonical compiler establishes
completeness. The acceptance change is exhibited by a concrete certificate
in the validation artifacts.

\section{New result: a smaller byte implementation}
\label{sec:byte-proof}

The incumbent recursive checker returns \texttt{None} when successful.
Its major-supply expression therefore uses
\texttt{c(m,g,k) or m*3}: after verification the expression has the value
$m*3$. The new byte profile makes every successful return truthy:
it returns \texttt{True} on a hypothesis, and otherwise a singleton list
containing an ordered pair. The pair's first component computes the
assumption supply; its second performs the recursive continuation.

Consequently the matching expression is
\texttt{c(m,g,k) and m*3}. If $m=0$, it still evaluates to zero and the
old arithmetic/introduction alternatives are selected. If $m\ne0$, it
evaluates to the same tripled major. A failed recursive verification still
raises; failure is never represented by a false return value that could
be mistaken for an axiom branch.

\begin{theorem}[Exact successful-execution correspondence]
The byte-profile checker and its 506-byte control accept the same finite
certificates, idealizing sufficient memory and recursion resources. They
consume the same suffix of the certificate and check the same scope and
context conditions in the same order.
\end{theorem}
\begin{proof}
Induct on the finite successful invocation tree in either direction.
A hypothesis changes the return value from \texttt{None} to
\texttt{True}, without reading a packet. Otherwise the singleton
comprehension reads one packet. Its ordered pair first computes the supply
expression. Successful major checks are truthy by induction, so the
Boolean expression produces exactly the old supply value. A DI assignment
to the goal occurs at this point, before the continuation.

Python evaluates the pair's second component only after its first.
Within the continuation call its arguments are evaluated left to right.
Thus source/template opening still sees the old floor, then the fresh
floor is assigned, and the final goal and new assumption openings see
that new floor. The local opener's closure observes the same floor
updates as in the control. The recursive invocation therefore receives
the same goal, context, and floor; induction applies.

Every successful nonhypothesis return is a list of exactly one pair, hence
truthy independently of the pair's contents. Conversely a completed pair
contains a completed recursive continuation, so the argument reverses.
No exception is caught by this transformation, and the driver emits
\texttt{ACCEPT} only after the root call completes. This proves the
claimed acceptance equivalence and consumption correspondence.
\end{proof}

The 503-byte profile retains the incumbent spelling \texttt{\~{}-k} for
$k-1$. The scope floor is always an integer, so these expressions agree.
A companion uses the ordinary spelling $k-1$ and measures 505 bytes;
it also beats the old 506-byte profile. The 505-byte companion is preferable
when source readability matters more than the two additional bytes.
Neither profile improves every coordinate of the incumbent; the token
and AST improvements come from the inferred-major profiles.

\section{A smaller complete native public-goal checker}
\label{sec:native-proof}

The separate native checker validates the Eq/$\D$ goal before interpreting
its proof. Its inherited opener carries a private mode $v$: initially
$v=1$, and $v=0$ throughout subsequent proof checking. The first scan
checks list arities, canonical integer-coordinate spellings, and scope.
The existing depth parity records the formula/term sort, so no independent
parser or grammar file is imported.

For a Boolean-valued predicate $P$, the guard
\[
            \neg v\lor P\qquad\hbox{is equivalent to}\qquad v\le P
            \quad(v\in\{0,1\}).
\]
The new implementation uses that ordered Boolean guard in the two
validation positions and replaces a join of two bits by bitwise union.
The required safety fact is not merely Boolean algebra: the right-hand
predicate is now evaluated even when $v=0$.

At a list branch, the left discriminator has already established an
actual list, so the arity predicate is total. At a scalar branch, the
scope comparisons have succeeded; the remaining public predicate first
tests exact integer type and then numerical parity and bounds, so it is
total whenever reached. The predicates are side-effect free. At the
fallback, both operands are bits ($d\bmod2$ and $v$), so logical and
bitwise union have the same zero/nonzero result.

\begin{theorem}[Native public-checker preservation]
The new native profile accepts exactly the finite JSON certificates
accepted by the pinned 656-byte native control, under the same resource
idealization. In particular it retains the complete public-goal grammar
and scope validation, and it presents the same full native theory.
\end{theorem}
\begin{proof}
The private mode invariant follows directly from the driver, default
argument, and reset before the first packet. The preceding totality and
Boolean identities prove equality of the initial validation outcome and
all subsequent opening outcomes. Every packet read, major recursion,
assumption supply, scope change, and continuation is otherwise unchanged.
Induction on a finite consumed execution establishes the equivalence.
\end{proof}

``Complete public-goal validation'' has its precise inherited meaning.
Unused certificate prefixes are allowed, and not every irrelevant
internal JSON value is lexically revalidated. The new implementation
does not relax those conditions; its entire input acceptance relation
equals its control's. Nor does this establish a unique spelling for all
dummy numerical slots. The abstract native language and the sorted
certificate interpretation specify their meaning.

\section{Measurements and reproducibility}
\label{sec:measurements}

All figures count an entire standalone acceptance module, including its
JSON input and success verdict. Each uses only the Python standard
library. The shared Python runtime and standard library are excluded on
both sides, exactly as in the pinned comparison corpus. No proof producer,
translator, normalizer, reference file, or test helper is imported by an
acceptance module.

The metric script removes actual docstrings, applies \texttt{ast.unparse},
adds one LF, and counts UTF-8 bytes, significant Python tokens, and AST
nodes. Comments, whitespace tokens, and end markers do not count as
significant tokens. The reference runtime is CPython 3.12.14. The archive
contains the measured sources, metrics in machine-readable form, and the
measurement script. Source is compiled before measurement is accepted;
parsing into an AST alone does not establish that Python accepts the
program.

\begin{center}\small
\begin{tabular}{lrrr}
\toprule
Complete module & Tokens & AST nodes & Bytes\\
\midrule
Token control & 256 & 283 & 594\\
Inferred-major token profile & 253 & 280 & 599\\
AST control & 260 & 278 & 602\\
Inferred-major AST profile & 257 & 275 & 607\\
Byte control & 267 & 303 & 506\\
Truthy-result byte profile & 278 & 310 & 503\\
Byte profile, ordinary decrement & 278 & 309 & 505\\
Native public-goal control & 333 & 379 & 656\\
Native public-goal profile & 327 & 375 & 643\\
\bottomrule
\end{tabular}
\end{center}

The identified compact coordinate envelope moves from
\[
      (256,278,506)\quad\hbox{to}\quad(253,275,503)
\]
in the order significant tokens, AST nodes, normalized bytes. The ordinary
decrement companion gives a 505-byte option. These are coordinatewise
attained bounds, obtained by different source programs. A single joint
minimum is not claimed. Likewise the audit establishes the listed
incumbents, not an exhaustive lower-bound theorem over all possible
implementations or every historical file in the repository.

The native profile improves its matching fully validated incumbent from
$(333,379,656)$ to $(327,375,643)$. Its larger size is not evidence that
it implements a weaker mathematical theory: it performs additional
input validation and accepts unrestricted native Eq/$\D$ public goals.
It should not be compared directly to a compact profile that assumes the
guarded public contract.

\subsection{Executable verification}

The delivered harness compiles each standalone source, remeasures it,
runs inherited positive and negative certificates, tests the inferred-major
expansion, and compares the exact-preservation refinements against their
controls. The suite contains real proof steps using computation,
introduction, elimination, equality, and scope allocation, as well as
malformed or unproved goals. A specific changed-major example demonstrates
that the new inference protocol is exercised.

The validation results are recorded in \texttt{evidence/validation.json}.
They supplement the universal proofs above. A finite test suite cannot
establish soundness, completeness, bi-interpretability, or a global source
minimum. Ordinary and optimized Python runs are included where relevant
to check that acceptance does not depend on assertions being enabled.

To reproduce the deliverable, run
\begin{lstlisting}
python3 reproduce.py
python3 reproduce.py --pdf
\end{lstlisting}
The first command performs the executable checks and measurements; the
second additionally compiles this source with pdfLaTeX. The provided PDF
is compiled from the provided LaTeX, with its pages rendered and inspected.
The source includes its appendices and checker listings directly and
requires no network access to compile.

\section{Assessment and remaining limits}

Under the user's stipulated implementation measure, the report makes a
strict reduction below each identified compact coordinate incumbent, and
a reduction of all three coordinates of one complete native public-goal
checker. The changes are small and exactly described. The inferred-major
protocol removes redundant certificate data; the other two refinements
preserve existing certificates. These are the new results.

The work does not establish an absolutely simplest axiomatization,
a global program-size lower bound, a single program simultaneously below
all three compact incumbents, or a new deletion of a primitive native
axiom. It does not extend the guarded public bi-interpretation to every
administrative raw object. It also does not treat the inherited Lean
development as a machine verification of the new Python source.

The ordinary theory is fully specified, the two interpretations and their
composites are included, and the new implementation arguments preserve
that mathematical content. That is a reproducible improvement under the
specified metric; its limited mathematical novelty should remain visible
alongside the numerical improvement.

\clearpage
\appendix
\section{Attribution and reading the inherited appendices}

The following nine sections reproduce the pinned ProveIt foundations,
with their original mathematical prose and internal cross-references.
Their statements described as ``new'' or ``the new calculus'' refer to
the historical packages from which they came, not to the present report.
They are included so that the interpretation and arithmetic dependencies
can be checked without retrieving another archive. The MIT No Attribution
license is supplied as \texttt{LICENSE-ProveIt.txt}; exact source paths and
Git blob hashes are in \texttt{evidence/source\_provenance.json}.

The contribution of the present report is confined to
Sections~\ref{sec:inferred}, \ref{sec:byte-proof}, and
\ref{sec:native-proof}, their measurements, and their executable artifacts.
The predecessor, equality characteristic, logical macros, guarded
translation, source admissibility, and scalar primitive-recursion
compiler are inherited. This distinction is also recorded in
\texttt{CLAIMS.json}.

% New mathematical result: complementary introduction choices.
% Requires amsmath/amssymb/amsthm and theorem, lemma, proposition environments.
% The inherited arithmetic-to-iPRA bridge must be supplied separately.

\section{Complementary introductions with one local assumption}
\label{sec:choice-introduction}

The old paired introduction retains both the zero instance and an
induction instance when its matrix is an equation. Both together are
unnecessary. A certificate may choose either one, provided the induction
choice is restricted to a literal equation. This changes neither the
formula language nor the meaning of a formula. The omitted combination
is recovered by a short nested introduction--elimination construction.

\subsection{Complete native syntax and rule inventory}

The scalar numerical language is
\[
 t::=x\mid0\mid St\mid I(n,a;u.h).
\]
Only $h$ binds the state variable $u$; $n$ and $a$ do not.
The update may contain fixed numerical parameters. The computation
axioms in every retained context are
\[
 I(0,a;u.h)=a,\qquad
 I(Sn,a;u.h)=h[I(n,a;u.h)/u].
 \tag{I0, IS}\label{eq:choice-computations}
\]
The argument also applies unchanged to the full fixed-description PR
language with its projection, composition, and recursion equations.
No numerical variable ranges over function descriptions, and no uniform
evaluator is postulated.

Native formulas and reusable context clauses are
\[
 A::=(s=t)\mid\mathsf D_{x,y}A,\qquad
 \langle A;y.B(y)\rangle,\quad y\notin\operatorname{FV}(A).
\]
Bound variables are renamed before capture-free substitution.
The two variables $x,y$ in a paired formula are distinct.
An ordinary assumption $F$ is the constant clause
$\langle H;v.F\rangle$, where $v$ is absent from $F$ and $H$ is the
fixed computation theorem $I(0,0;u.0)=0$. Thus ordinary hypothesis
use is a clause Use preceded by this computation proof.

There are computation axioms and exactly the following inference
patterns:
\[
\frac{\Gamma\vdash A}{\Gamma\vdash B(t)}
\quad\mathsf{Use}
\quad(\langle A;y.B(y)\rangle\in\Gamma),
\]
\[
\frac{\Gamma\vdash\mathsf D_{x,y}B(x,y)\qquad
      \Gamma,\langle B(a,0);y.B(a,y)\rangle\vdash F}
     {\Gamma\vdash F}
\quad\mathsf{DE},
\qquad
\frac{\Gamma\vdash s=t\qquad\Gamma\vdash F(s)}
     {\Gamma\vdash F(t)}
\quad\mathsf{Tr},
\]
\[
\frac{\Gamma,B(t,\eta_\varepsilon(z))\vdash B(t,Sz)}
     {\Gamma\vdash\mathsf D_{x,y}B(x,y)}
\quad\mathsf{DI}_{\varepsilon},
\qquad
\eta_0(z)=0,\quad\eta_1(z)=z.
\tag{Choice DI}\label{eq:choice-di}
\]
Here $\varepsilon$ is a certificate choice in $\{0,1\}$, and
$\varepsilon=1$ is permitted only when $B$ has a literal equation
root. An equation beneath a binder is not a literal equation.
Call $\varepsilon=0$ the \emph{case choice} and $\varepsilon=1$
the \emph{progress choice}. Both retain exactly one ordinary
formula assumption. Compound matrices have only the case choice.

In DI the eigenparameter $z$ is fresh for $\Gamma$, the witness
$t$, and every fixed parameter of $B$. In DE the eigenparameter
$a$ is fresh for $\Gamma$, the conclusion $F$, and every fixed
parameter of $B$. In transport, $F(v)$ is any fixed native
formula template. Displayed local assumptions and clauses are
discharged. There are no primitive reflexivity, coverage,
separation, ordinary logical, or additional induction rules.

We call this calculus $\mathcal C$. The original comparison
calculus $\mathcal N$ has identical syntax, computation,
Use, DE, and transport, but its introduction is
\[
\frac{\Gamma,B(t,0),[B(t,z)]_{\mathrm{at}}\vdash B(t,Sz)}
     {\Gamma\vdash\mathsf D_{x,y}B(x,y)}.
\tag{Old DI}\label{eq:choice-old-di}
\]
The bracket supplies $B(t,z)$ exactly at a literal equation.

\subsection{Exact equivalence and the unused-witness construction}

\begin{lemma}[Structural admissibility]
\label{lem:choice-structural}
Weakening, exchange, contraction, renaming, ordinary-assumption cut,
and capture-free numerical substitution are admissible in $\mathcal C$.
\end{lemma}
\begin{proof}
Induct on derivations. Rename local eigenparameters before inserting
proofs or terms. Numerical substitution preserves a formula's root
constructor, so an allowed progress choice remains allowed.
The selected choice does not change under substitution or cut.
For ordinary-assumption cut, replace each use of its constant clause
by the supplied proof. Computation schemas are closed under
capture-free numerical substitution. These facts verify every case.
\end{proof}

\begin{theorem}[Identity-language equivalence]
\label{thm:choice-equivalence}
For every retained clause context $\Gamma$ and native formula $F$,
\[
 \Gamma\vdash_{\mathcal C}F
 \quad\Longleftrightarrow\quad
 \Gamma\vdash_{\mathcal N}F.
\]
The transformations are effective on finite derivations and leave
all numerical and formula syntax unchanged.
\end{theorem}
\begin{proof}
Each new introduction is an instance of the old one after weakening
by the omitted local assumption. For the case choice at an equation,
add the induction instance. For the progress choice, add the zero
instance. Compound case introduction is already the old rule.
Thus every new derivation translates to an old one.

For the reverse direction, transform an old derivation recursively.
Only its atomic DI case needs a construction. Suppose its matrix is
the equation $E(x,y)$ and its transformed premise is
\[
 \pi:\quad\Gamma,E(t,0),E(t,z)\vdash E(t,Sz).
\]
The outer new inference will be case DI with witness $t$. It is
therefore enough to prove $E(t,Sz)$ under
$\Gamma,E(t,0)$.

Rename the eigenvariable in $\pi$ to a fresh $v$. Choose a new
variable $u$ absent from $E(t,y)$ and all retained data. Under
$\Gamma,E(t,0)$, progress DI with witness zero introduces
\[
             \mathsf D_{u,y}E(t,y)
\]
from exactly the renamed premise $\pi[v/z]$.
The bound coordinate $u$ is deliberately unused.

Apply DE to this paired formula. Because $u$ is unused, the opened
clause is exactly
\[
               \langle E(t,0);y.E(t,y)\rangle;
\]
it contains no new existential witness. The retained ordinary
assumption $E(t,0)$ supplies its antecedent. Use at $Sz$ gives
$E(t,Sz)$, and DE closes. Finally the outer case DI introduces
the original $\mathsf D_{x,y}E(x,y)$ with witness $t$.
The original premise proof $\pi$ has been inserted once.

An old compound introduction already is the new case choice.
All other rule cases are unchanged. This completes the induction
on the original derivation.
\end{proof}

The unused inner coordinate is essential. Introducing
$\mathsf D_{x,y}E(x,y)$ internally and then eliminating it would
give an arbitrary fresh witness, not necessarily the original
$t$. The construction instead keeps $t$ as a fixed parameter
of $E(t,y)$ and binds only the unused coordinate $u$.

\begin{proposition}[Linear bound on inference nodes]
\label{prop:choice-node-bound}
Count each computation axiom, Use, DI, DE, and transport as one
inference node, expanding ordinary-assumption use as a computation
of $H$ followed by Use. If an old derivation has $N$ nodes and
$r$ atomic DI nodes, the displayed transformation has exactly
\[
                   N+5r\leq 6N
\]
nodes.
\end{proposition}
\begin{proof}
Each atomic old DI node is replaced by six nodes: outer case DI,
inner progress DI, DE, Use of the opened clause, Use of the
retained base assumption, and the computation of $H$ needed for
that last ordinary-assumption use. Its old premise subtree
occurs once. All other inference nodes are copied once.
\end{proof}

This counts inference-tree nodes. It is not a bound on serialized
JSON bytes, formula sizes, checker runtime, or the number of
consumed packet cells when hypothesis closure is implicit.
In the formula-context Q/D presentation, the corresponding
replacement has outer case DI, inner progress DI, DE, QE,
and two Hyp nodes, giving the same six-node count.

\subsection{Induction and numerical cases separate cleanly}

\begin{lemma}[Contextual equation induction]
\label{lem:choice-ei}
The progress choice, DE, and Use derive
\[
\frac{\Gamma\vdash E(0)\qquad
      \Gamma,E(z)\vdash E(Sz)}
     {\Gamma\vdash E(t)}
\tag{EI}\label{eq:choice-ei}
\]
for every equation $E$, with the usual freshness condition.
\end{lemma}
\begin{proof}
Choose an unused witness coordinate $u$. Progress DI introduces
$\mathsf D_{u,y}E(y)$ from the given step. DE opens the clause
$\langle E(0);y.E(y)\rangle$. Supply its antecedent using the given
base proof and use the clause at $t$. Close DE.
\end{proof}

\begin{lemma}[Cases, equality, and binder congruence]
\label{lem:choice-preinduction}
The case choice derives, for arbitrary native formulas,
\[
\frac{\Gamma\vdash F(0)\qquad\Gamma\vdash F(Sz)}
     {\Gamma\vdash F(t)}
\tag{NC}\label{eq:choice-nc}
\]
with the usual freshness condition. Equality reflexivity,
symmetry, transitivity, numerical congruence, and uniform
matrix interderivability under $\mathsf D$ are derivable
without using the progress choice.
\end{lemma}
\begin{proof}
For cases, use case DI with an unused witness coordinate and
matrix $F(y)$, ignoring its retained base in the successor
premise. DE opens $\langle F(0);y.F(y)\rangle$; the given zero
proof and Use at $t$ finish.

For reflexivity put $r=I(0,t;u.0)$. Use the computation equation
$r=t$ as both transport premises with motive $v=t$ to obtain
$t=t$. The usual explicit equality motives give symmetry,
transitivity, and numerical congruence.

For binder congruence, open the old paired formula at a fresh
witness $a$ and introduce the new one by case DI at $a$.
Convert its retained base to the old base, use the opened
old clause at the successor, and convert back. Ordinary cut
closes the construction, and the reverse conversion gives
the converse. No induction is used.
\end{proof}

The primitive progress choice is still restricted to equations.
Allowing it at arbitrary compound matrices would, by the
argument for \eqref{eq:choice-ei}, supply induction for
arbitrary formulas and change the intended theory.

\subsection{The entire ordinary logical interface}
\label{subsec:choice-logic}

Here is a predecessor-free reconstruction of the ordinary logical
operations. It makes explicit that the case choice supplies the
logical interface independently of progress induction.

Let $C(v,a,b):=I(v,a;u.b)$ with $u$ unused. Computation gives
$C(0,a,b)=a$ and $C(Sv,a,b)=b$. Define the gate on fixed finite
formula syntax by
\[
 G_v(s=t):=(C(v,0,s)=C(v,0,t)),\qquad
 G_v(\mathsf D_{x,y}B):=\mathsf D_{x,y}G_vB,
\]
renaming bound variables away from the control. Equality,
binder congruence, and external structural induction give
\[
 \vdash G_0A,\qquad G_{Sv}A\dashv\vdash A,\qquad
 A\dashv\vdash\mathsf D_{u,y}G_yA
 \quad(u,y\notin\operatorname{FV}(A)).
\tag{Gate and padding}\label{eq:choice-gate}
\]
For padding, case DI proves the forward direction. DE, the
zero-gate theorem, and Use at $1$ prove the reverse direction.

Define a selector between two arbitrary formulas by
\[
 M_v(s=t,u=w):=(C(v,s,u)=C(v,t,w))
\]
and, with aligned fresh paired variables, by
\[
\begin{aligned}
 M_v(\mathsf D_{x,y}A,\mathsf D_{x,y}B)
   &:=\mathsf D_{x,y}M_v(A,B),\\
 M_v(E,\mathsf D_{x,y}B)
   &:=\mathsf D_{x,y}M_v(G_yE,B),\\
 M_v(\mathsf D_{x,y}A,E)
   &:=\mathsf D_{x,y}M_v(A,G_yE),
\end{aligned}
\]
where $E$ is an equation. Every call decreases maximum input
binder depth. Padding and congruence prove
$M_0(A,B)\dashv\vdash A$ and
$M_{Sv}(A,B)\dashv\vdash B$.

Put $\delta(v)=C(v,1,0)$ and $p(z)=I(z,0;u.\delta(u))$.
Only the fixed computations $p(1)=1$ and $p(2)=0$ are needed.
The full logical macro inventory is
\[
\begin{aligned}
 A\Rightarrow_*B&:=\mathsf D_{u,z}M_z(A,B),\\
 A\wedge_*B&:=\mathsf D_{u,z}G_zM_{p(z)}(A,B),\\
 \exists_*x\,A(x)&:=\mathsf D_{x,y}G_yA(x),\\
 A\vee_*B&:=\exists_*x\,M_x(A,B),\\
 \forall_*x\,A(x)&:=
        A(0)\wedge_*\mathsf D_{u,z}G_zA(z),\\
 \bot_*&:=(0=1).
\end{aligned}
\tag{Logical macros}\label{eq:choice-macros}
\]
Newly bound variables are fresh; $u$ is unused.

Implication introduction obtains its assumption from the zero
selector and supplies the consequent at successors, then uses
case DI. Elimination opens its clause and uses it at $1$.
Existential introduction uses its witness and the successor
gate; elimination opens at a fresh witness, supplies the
zero-gate theorem, and uses the clause at $1$.

Disjunction introduction uses witnesses $0$ and $1$.
For elimination, discharge the two branch proofs as
$A\Rightarrow_*F$ and $B\Rightarrow_*F$. NC applied to
$M_z(A,B)\Rightarrow_*F$ gives this implication at every
control: the two endpoint cases follow from the branch
implications. Existential elimination provides $M_a(A,B)$
and the implication gives $F$.

Given $A$ and $B$, selector endpoints and NC give $M_v(A,B)$
at an arbitrary control. Substitute $v=p(Sz)$ and use the
successor gate for conjunction introduction. For projections,
open its zero-gated clause and specialize at $1$ and $2$;
the two fixed computations give $B$ and $A$, respectively.

These rules prove existential numerical coverage:
\[
     (x=0)\vee_*\exists_*v\,(x=Sv).
\tag{Coverage}\label{eq:choice-coverage}
\]
Use NC. Zero takes the left injection of reflexivity, and a
successor takes the right injection with the evident witness.
Coverage gives equality cases even when the retained context
contains $x$; its elimination adds $x=0$ or $x=Sv$ rather
than silently changing that retained context.

Universal introduction follows from a fresh-parameter proof
by substituting zero and successors and forming the two
conjuncts. For elimination, take its $A(0)$ conjunct and
open the paired conjunct. Its base is a zero-gate theorem.
Coverage of the target chooses the zero instance from
$A(0)$ or a successor instance from the opened clause
and gate removal.

Finally, under $0=1$, computation and transport give
$s=C(0,s,t)=C(1,s,t)=t$ for arbitrary numerical terms.
Every equation follows, and external induction on formula
syntax with case DI gives every native formula. This
supplies falsity elimination. No excluded middle for
quantified formulas has entered the construction.

\subsection{Whole-domain realizations and both formula composites}

For a standard first-order reading, let $T_{\rm eq}$ have
intuitionistic equality logic, the same numerical signature
and computation equations, numerical coverage, separation
$(0=1)\to\bot$, and contextual equation induction.
Its interpretation of the native syntax is
\[
 (s=t)^\circ=(s=t),\qquad
 (\mathsf D_{x,y}B)^\circ
   =\exists x\bigl(B^\circ(x,0)\to\forall y B^\circ(x,y)\bigr),
\]
\[
 \langle A;y.B\rangle^\circ
       =A^\circ\to\forall y B^\circ.
\]
The domain is the entire numerical domain and interpreted
equality is actual equality.

Case DI is sound in this realization by numerical coverage.
Progress DI is sound by contextual equation induction under
the conditional base. Use, DE, transport, and computation
have their displayed ordinary meanings. Conversely, the
native macros derive every ordinary logical rule,
\eqref{eq:choice-ei} derives equation induction, and
\eqref{eq:choice-coverage} derives coverage. Separation
translates to the identity implication
$(0=1)\Rightarrow_*(0=1)$.

For the essential paired formula composite,
\[
 \mathsf D_{x,y}B
 \dashv\vdash
 \exists_*x\bigl(B(x,0)\Rightarrow_*\forall_*y\,B(x,y)\bigr).
\tag{Paired composite}\label{eq:choice-paired-composite}
\]
Forwards, DE opens its clause at a fresh $a$; under $B(a,0)$
that clause supplies $B(a,y)$ at a fresh arbitrary $y$.
Introduce the universal, implication, and existential
macros and close DE. Backwards, eliminate the existential
at a fresh $a$. Introduce the paired formula with witness
$a$ by case DI: the retained base, supplied implication,
and universal elimination at $Sz$ give its successor premise.

The ordinary interpretations of the six logical macros
have their ordinary round trips. For conjunction, its true
base and the instances $1,2$ recover both operands;
the converse uses selector cases. The universal has all
successor instances in its paired conjunct and zero in
its first conjunct; coverage gives exactly the ordinary
universal. Implication, disjunction, and existential
quantification use their selector or gate endpoints.
Falsity uses the explicit ordinary separation bridge.
External induction on finite formulas, together with
\eqref{eq:choice-paired-composite}, gives both formula
composites. Clauses correspond in both directions to
$A\Rightarrow_*\forall_*yB$ by the derived logical rules.

The numerical composites are literal identities with
internally provable graph $x=y$. No standard-numeral
restriction, retract, quotient, or equality reflection
assumption is being inserted at this stage. The separately
supplied arithmetic bridge identifies $T_{\rm eq}$ with
full iPRA; the present change preserves that bridge
unchanged because Theorem~\ref{thm:choice-equivalence}
preserves every native sequent.

\subsection{What the theorem does and does not claim}

The new inference interface retains one ordinary assumption
at every paired introduction. It preserves the entire
consequence relation, not just equations or closed
theorems. Consequently the established full-iPRA
interpretations and equiconsistency transfer unchanged.

This result is different from retaining only progress
introduction and strengthening closure elimination.
Here the original elimination rule is retained, and
the case and progress introductions coexist. It is also
different from allowing arbitrary compound progress,
which would supply stronger induction.

A complete-checker comparison must still count the
certificate choice, its equation-root restriction,
substitution, context handling, computation, and the
public input/output interface. The mathematical theorem
alone is not a byte-count or minimality theorem.
The inference-node bound is distinct from source-size
measurements and from serialized certificate size.
The proof is a written universal proof; finite regression
tests do not constitute a Lean/Rocq formalization of it.


\section{Arithmetic strength and both ordinary interpretation composites}
\label{sec:ordinary-bridge}

We spell out the arithmetic bridge to avoid identifying standard numerical
agreement with bi-interpretability. First elaborate every fixed iterator
into an ordinary function symbol. At stage zero the symbols are $0,S$.
Given a fixed earlier-stage term $h(u,\vec p)$, introduce
$i_h(n,a,\vec p)$ with the two equations
\[
i_h(0,a,\vec p)=a,\qquad
i_h(Sn,a,\vec p)=h(i_h(n,a,\vec p),\vec p).
\]
Only finite descriptions occur; the domain contains no code for a
universal evaluator. Formal parameters are distinct when a description
is declared, while actual arguments can coincide.

Let $T_I$ be ordinary intuitionistic first-order logic with actual
equality and genuine falsity, these staged equations, contextual
induction for literal equations, and
\[
\forall x\,(x=0\vee\exists y\,x=Sy),
\qquad (0=S0)\to\bot.
\]
The second sentence is the falsity bridge. It must be stated: positive
numerical equations by themselves admit a singleton structure. Native
explosion from $0=S0$ and a translation of ordinary falsity to that
equation do not authorize silently dropping the bridge from $T_I$.

The native-to-ordinary translation is
\[
(s=t)^\circ=(s=t),\qquad
(\mathsf D_{x,y}B)^\circ=
\exists x\bigl(B^\circ(x,0)\to\forall y B^\circ(x,y)\bigr).
\]
A clause $\langle A;y.B\rangle$ becomes $A^\circ\to\forall yB^\circ$.
For the auxiliary closure used by the compact checkers,
$(\mathsf Q_yB)^\circ=B^\circ(0)\to\forall yB^\circ$.
The interpretation domain is all elements, defined by $x=x$;
equality is actual equality.

Case DI is sound because numerical coverage combines its retained base
with its successor proof. Progress DI is sound by literal-equation
induction under a temporarily assumed base. Use, DE and transport are
ordinary logical inferences under the translations. Conversely the
derived logical operations, equation induction, cases and explosion in
Section~\ref{sec:choice-introduction} recover the rules and axioms of
$T_I$. An auxiliary closure is interderivable with
$\mathsf D_{v,y}B(y)$ for unused $v$: in one direction use case DI with
the closure and its base; in the other eliminate the paired formula
and return its supplied closure. Thus formulas and context clauses have
the stated meanings in both presentations.

\subsection{The arithmetic dependency order}

The scalar update has no recursion-counter argument. We do not conceal
one in an encoding. The inherited internal proofs, reproduced in the
appendices, establish the required arithmetic in this order:
\begin{enumerate}
\item Using equation induction and numerical cases, construct a fixed
scalar predecessor term and prove $P0=0$ and $P(Sn)=n$.
Existential coverage then gives $x=0\vee x=S(Px)$: in the successor
case $x=Sy$, apply $P$ to obtain $Px=y$ and substitute.
\item Construct truncated subtraction and a numerical characteristic
of actual equality. Prove its reflection on the entire original domain.
\item Compile each fixed quantifier-free formula into its characteristic
equation and recover every contextual quantifier-free induction instance.
\item With that induction available, construct pairing and projections
by a bounded square-root search and compile every fixed general
primitive-recursive description into a scalar iterator term.
\end{enumerate}
No pairing or faithful equality detector is assumed to obtain
predecessor. No appeal to standard-model correctness replaces an
internal equality. Every induction fixes its passive parameters;
diagonal substitution occurs only into a previously proved identity.

For a target recursion
\[
f(0,\vec p)=g(\vec p),\qquad
f(Sn,\vec p)=h(n,f(n,\vec p),\vec p),
\]
let $\pi,L,H$ be the proved pairing and projections. The compiler uses
\[
V(n,\vec p)=I\bigl(n,\pi(0,g^*(\vec p));
u.\pi(SL(u),h^*(L(u),H(u),\vec p))\bigr),
\qquad f^*(n,\vec p)=H(V(n,\vec p)).
\]
Equation induction proves $L(V(n,\vec p))=n$; the two projection
identities then prove the target defining equations. The construction
is external recursion on each finite PR description, not an internal
quantification over arbitrary functions.

\begin{theorem}[Whole-domain bi-interpretability and equiconsistency]
The ordinary realization $T_I$ and full first-order iPRA are
bi-interpretable. In both theories $x=y$ defines the isomorphism from
the original domain to its double interpretation. The native Choice-DI
calculus is therefore a faithful proof presentation of a theory
bi-interpretable and equiconsistent with iPRA.
\end{theorem}
\begin{proof}
Interpret an iterator in iPRA by primitive recursion that ignores its
counter argument in the update. iPRA proves its equations, literal
equation induction, coverage and separation. In the reverse direction
use the scalar compiler. The predecessor and equality-reflection proofs
supply the successor axioms and every quantifier-free induction instance.
All ordinary logical rules are preserved.

For two iterators with the same seed and provably equal updates,
equation induction on the equality of their values proves uniqueness at
every counter, with all parameters fixed. The corresponding argument
applies to ordinary primitive recursion. Induct externally on a finite
description: its double translation satisfies its original defining
equations, so internal uniqueness identifies the two functions. This
also proves coherence when substitution changes the staged syntactic
description of an update.

The domain and equality are unchanged; hence the identity graph is
internally total, bijective and preserves every interpreted operation.
The displayed ordinary meanings of the logical macros and the paired
binder give both formula composites by external induction on finite
formula syntax. Clauses are similarly equivalent to their conditional
universal meanings. The explicit falsity bridge preserves contradiction
in both directions. Effective translations of finite proofs give
equiconsistency, without a claim that either theory proves its own
consistency.
\end{proof}


% Adapted with attribution from the pinned unparenthesized one-hole archive.
\section{Internal scalar predecessor bootstrap}
\label{sec:hybrid-bootstrap}

This scalar bootstrap is inherited from the supplied scalar development~\cite{erasure} and is reproduced to keep the scalar checker self-contained. It is not a new contribution of this paper. The proof uses a factor map fixing odd states and a saturated-iteration construction of predecessor.

The available arithmetic terms are $0$, $S$, and
\[
 I(n,a;u.h),\qquad
 I(0,a;u.h)=a,\qquad
 I(Sn,a;u.h)=h[I(n,a;u.h)/u].
\]
The update binds only $u$.  All other variables in it are fixed numerical
parameters.  We use equality reasoning, the already reconstructed
zero/successor cases, and induction on a single literal equation.
Every induction below keeps its other parameters fixed.  Predecessor,
successor injectivity, pairing, and induction on arbitrary
quantifier-free formulas are not assumptions of the argument.

\subsection{Elementary terms and two finite cycles}

Write $1=S0$ and use larger numerals as successor abbreviations.  Define
\[
 C(n,a,b)=I(n,a;u.b),\qquad
 Z(n)=C(n,1,0),\qquad
 a+b=I(b,a;u.Su).
\]
The update variable of $C$ is absent from its body.  Consequently
\[
 C(0,a,b)=a,\qquad C(Sn,a,b)=b.
\]
Numerical cases show that $Z(n)$ is $0$ or $1$.
The usual commutative-monoid equations for addition follow by literal
equation induction.  For example, $0+b=b$ and $Sa+b=S(a+b)$ are
proved by induction on $b$; these yield commutativity by a further
equation induction, and associativity follows by induction on the last
summand.  We use these equations to reorder finite sums.  The notation
$2a$ means $a+a$ and $3a$ means $a+a+a$; no multiplication operation
is needed for this bootstrap.

Put
\[
 N(n)=I(n,0;u.Z(u)),\qquad
 j(u)=Z(u)+2N(u),\qquad
 M(n)=I(n,0;u.j(u)).
\]
The range of $N$ is contained in $\{0,1\}$: its zero value is $0$,
and its successor value is a value of $Z$.  This uses numerical cases
on its counter, not induction on a disjunction.  Hence
\begin{equation}
 N(n)+Z(N(n))=1,\qquad
 Z(Z(N(n)))=N(n),\qquad N(n+2)=N(n).
 \label{eq:hybrid-parity}
\end{equation}
Here the last equation follows from two successor computations and
the preceding equation.  Literal-equation inductions give
\begin{equation}
 N(2n)=0,\qquad N(2n+1)=1,\qquad N(3n)=N(n).
 \label{eq:hybrid-parity-scaling}
\end{equation}
For the first equation the successor step adds two to the argument of
$N$; the second follows by one computation.  In the third, increasing
$n$ adds three to the left argument and therefore applies $Z$ once to
both sides of the induction equation.

Finite computations now give
\[
 j(0)=1,\qquad j(1)=2,\qquad j(2)=0.
\]
The proof below needs just these three values and
$M(Sn)=j(M(n))$; it does not assume a division algorithm or a theorem
characterizing $M$ as remainder modulo three.

\subsection{A ramp with a global scaling identity}

Define
\[
 G(n)=I(n,0;u.u+M(u)+1).
\]
The central invariant is
\begin{equation}
 M(G(n))=N(n).
 \label{eq:hybrid-ramp-invariant}
\end{equation}
Its zero case is immediate.  For its successor step, retain the
single induction equation and split the Boolean value $N(n)$.
If it is $0$, the next $G$ state is $G(n)+1$; its $M$ value is
$j(0)=1$, equal to $N(Sn)$.  If it is $1$, the next $G$ state
is $G(n)+2$; two computations of $M$ change its value from $1$
to $2$ and then to $0$, again equal to $N(Sn)$.  This proves
the invariant by literal-equation induction.

It follows that
\[
 G(Sn)=G(n)+N(n)+1,\qquad G(n+2)=G(n)+3.
\]
Equation induction therefore yields
\begin{equation}
 G(2n)=3n.
 \label{eq:hybrid-ramp-double}
\end{equation}
A separate literal-equation induction proves
\begin{equation}
 G(3n)=3G(n)+N(n).
 \label{eq:hybrid-ramp-triple}
\end{equation}
For completeness, its successor calculation is
\begin{align*}
 G(3Sn)
 &=G(3n)+N(3n)+4
   =3G(n)+2N(n)+4,\\
 3G(Sn)+N(Sn)
 &=3G(n)+3N(n)+3+Z(N(n)).
\end{align*}
The expressions agree in both Boolean cases for $N(n)$.
The base is $0=0$.  In particular,
\eqref{eq:hybrid-ramp-triple} was proved internally; it was not
inferred from an intended floor-function interpretation of $G$.

Set
\[
 H(u)=C(N(u),G(u),u).
\]
The crucial choice is to retain $u$ when its parity is $1$.
The preceding equations give
\begin{equation}
 H(2u)=3u,\qquad H(1)=1,\qquad H(3u)=3H(u).
 \label{eq:hybrid-factor-map}
\end{equation}
To check the last identity, use $N(3u)=N(u)$ and split the Boolean
$N(u)$.  At $0$, both sides are $3G(u)$ by
\eqref{eq:hybrid-ramp-triple}; at $1$, both are $3u$.
This argument does not require that an arbitrary number has already
been represented as $2v$ or $2v+1$.

\subsection{A comparator from scalar iteration}

For every fixed update $h$, equation induction on $n$, keeping $m,a$
and all parameters fixed, proves the semigroup identity
\[
 I(m+n,a;h)=I(n,I(m,a;h);h).
\]
The base is the zero computation and the step applies $h$ to the
induction equation. This identity will also be used in the later
pairing construction.


For a fixed update $h$, equation induction on $y$, with its initializer
$a$ and all parameters of $h$ fixed, proves the rotation identity
\begin{equation}
 I(Sy,a;h)=I(y,h(a);h).
 \label{eq:hybrid-rotation}
\end{equation}
Both sides have base $h(a)$, and their successor equations apply the
same $h$ to the previous equality.  Here and below $I(y,a;h)$ is
an abbreviation for one fixed finite scalar term, not quantification
over functions.

By \eqref{eq:hybrid-factor-map}, another equation induction gives
\begin{equation}
 I(y,3a;H)=3I(y,a;H).
 \label{eq:hybrid-iterator-scaling}
\end{equation}
Its base is $3a=3a$.  In the step the fixed-parameter induction
equation rewrites the new left side as
\[
 H(3I(y,a;H))=3H(I(y,a;H)).
\]
The initializer parameter stays $a$ throughout this induction.
Equation induction also gives $I(y,1;H)=1$, since $H(1)=1$.

Define
\[
 E(x)=I(x,1;u.2u),\qquad
 L(x,y)=N\bigl(I(y,E(x);H)\bigr).
\]
Then the following entire recursion specification is provable:
\begin{equation}
 L(0,y)=1,\qquad L(Sx,0)=0,\qquad L(Sx,Sy)=L(x,y).
 \label{eq:hybrid-comparator}
\end{equation}
The first equation uses $E(0)=1$ and the fixed-point identity for
$H$.  The second uses $E(Sx)=2E(x)$ and $N(2a)=0$.
For the third,
\begin{align*}
 L(Sx,Sy)
 &=N\bigl(I(Sy,2E(x);H)\bigr)\\
 &=N\bigl(I(y,H(2E(x));H)\bigr)\\
 &=N\bigl(I(y,3E(x);H)\bigr)\\
 &=N\bigl(3I(y,E(x);H)\bigr)\\
 &=N\bigl(I(y,E(x);H)\bigr)=L(x,y).
\end{align*}
The successive steps use rotation, $H(2a)=3a$, iterator scaling,
and $N(3a)=N(a)$.  Thus $L$ has the comparison equations at all
arguments internally.  Its standard-model meaning as the indicator
of $x\le y$ is motivation, not a premise of the following argument.

\subsection{Saturated iteration gives predecessor}

We now need only the three equations
\eqref{eq:hybrid-comparator}.  Define
\[
 B_m(u)=C(L(Su,m),u,Su),\qquad
 W(t,m)=I(t,0;u.B_m(u)).
\]
The subscript in $B_m$ records a fixed parameter; $B_m$ is an open
term abbreviation, not a function variable.

Numerical cases on $e$ prove the selector equation
\[
 C(e,Sa,Sb)=S C(e,a,b).
\]
Equivalently one may use equation induction on $e$, whose induction
hypothesis is unused.  Therefore
\begin{align}
 B_{Sm}(Su)
 &=C(L(SSu,Sm),Su,SSu)\notag\\
 &=C(L(Su,m),Su,SSu)=S B_m(u),
 \label{eq:hybrid-conjugacy}\\
 B_{Sm}(0)&=C(L(S0,Sm),0,S0)=C(L(0,m),0,S0)=S0.
 \label{eq:hybrid-conjugacy-base}
\end{align}

Literal-equation induction on $t$, with $m$ fixed, proves
\begin{equation}
 W(St,Sm)=S W(t,m).
 \label{eq:hybrid-saturated-shift}
\end{equation}
The base is \eqref{eq:hybrid-conjugacy-base}.  The successor step is
\begin{align*}
 W(SSt,Sm)
 &=B_{Sm}(W(St,Sm))\\
 &=B_{Sm}(S W(t,m))\\
 &=S B_m(W(t,m))=S W(St,m).
\end{align*}
Only the second line uses the current induction hypothesis, at the
same fixed $m$.

Two subsequent equation inductions give
\begin{equation}
 L(Sn,n)=0,\qquad W(n,n)=n.
 \label{eq:hybrid-diagonals}
\end{equation}
The first uses the comparator recurrence.  The second uses
\eqref{eq:hybrid-saturated-shift} in its successor step.  The
substitution $t=m=n$ is made into that already proved universally
parameterized equation; no parameter of an active induction hypothesis
is changed.  Consequently
\begin{equation}
 W(Sn,n)=B_n(n)=C(L(Sn,n),n,Sn)=n.
 \label{eq:hybrid-extra-step}
\end{equation}

Finally define the finite scalar term
\begin{equation}
 P(m)=I\bigl(m,0;u.C(L(SSu,m),u,Su)\bigr).
 \label{eq:hybrid-predecessor-term}
\end{equation}
Its zero equation is immediate.  To justify its successor equation
without assuming a separate extensionality rule for iterator bodies,
put
\[
 V(t,m)=I\bigl(t,0;u.C(L(SSu,Sm),u,Su)\bigr).
\]
The comparator recurrence identifies this update pointwise with
$B_m(u)$.  Equation induction on $t$, with $m$ fixed, gives
$V(t,m)=W(t,m)$: their base values are $0$, and the successor step
uses the pointwise update equation and then the induction equation.
Now \eqref{eq:hybrid-extra-step} yields
\[
 \boxed{\qquad P(0)=0,\qquad P(Sm)=V(Sm,m)=W(Sm,m)=m.\qquad}
\]

\subsection{Dependency boundary}

All terms in this section are built from $0,S,I$ by finite substitution.
All inductions have a single equation as their invariant.  Numerical
cases and equality transport are the already reconstructed logical
resources.  Successor injectivity may now be obtained by applying $P$
to an assumed equation $Sx=Sy$.  Equality tests, reflection of
quantifier-free formulas, pairing, and the recovery of general
primitive recursion can then be developed in the established order.
They were not used in order to manufacture predecessor.

This argument is a written internal derivation in the stated theory.
The finite implementation checks do not constitute a proof-assistant
formalization of this entire metatheorem.


% Standalone section.  Requires only amsmath and amssymb in the parent document.
% No theorem environments, macros, or bibliography keys are assumed.
% All arithmetic claims in this section are proved here.

\section{From literal-equation induction to full intuitionistic PRA}
\label{sec:eic-completion}\label{app:arithmetic}

This section supplies an unconditional completion theorem. It is adapted
with attribution from the current forward-left draft~\cite{forward}. In the
scalar application used here, every function occurring in the proof is
already a scalar term: predecessor is supplied by
Appendix~\ref{sec:hybrid-bootstrap}, the selector is
$C(n,a,b)=I(n,a;u.b)$, addition is $x+y=I(y,x;u.Su)$, and
truncated subtraction is $d(x,y)=I(y,x;u.Pu)$. All other functions below
are finite compositions of these. Thus this proof does not presuppose
that arbitrary primitive recursion has already been compiled. The QFI
conclusion first applies to the scalar language; the general PR compiler
comes afterward.  Literal-equation
induction, together with the primitive-recursive defining equations,
predecessor coverage, and separation of $0$ and $1$, proves induction for every
quantifier-free formula.  In particular, the proof does not assume a faithful
numerical equality test, stability or decidability of equality, successor
uniqueness, or any commutation law for an iterator.  We construct the equality
test and prove its reflection property for arbitrary elements of the
numerical domain.

\subsection{The exact hypotheses}

Work in ordinary intuitionistic first-order logic with equality and genuine
falsity $\bot$.  The language has one numerical sort and a symbol for each
fixed, finite primitive-recursive description.  Distinct descriptions are
not identified merely because they compute the same standard function.
Let $\mathsf{PRE}$ contain the defining equations for zero, successor,
projections, composition, and primitive recursion.  In particular, choose
descriptions $P$ and $C$ with equations
\begin{align}
 P(0)&=0, & P(Sx)&=x, \label{eq:eic-predecessor}\\
 C(0,a,b)&=a, & C(Sx,a,b)&=b. \label{eq:eic-selector}
\end{align}
Write $1=S0$.  In addition to $\mathsf{PRE}$, assume
\begin{align}
 \mathsf{Cov}:&\qquad
       \forall x\,\bigl(x=0\lor x=S(Px)\bigr),
       \label{eq:eic-coverage}\\
 \mathsf{Sep}:&\qquad 0=1\longrightarrow\bot,
       \label{eq:eic-separation}
\end{align}
and induction for each \emph{literal equation} $E(n,\vec p)$:
\begin{equation}
 \forall\vec p\,
 \left[
   \left(E(0,\vec p)\land
       \forall n\bigl(E(n,\vec p)\longrightarrow E(Sn,\vec p)\bigr)
   \right)
   \longrightarrow\forall n\,E(n,\vec p)
 \right].
 \label{eq:eic-induction}
\end{equation}
Here the schematic expression $E$ is a single equality of numerical
terms; the displayed universal quantifiers range only over numbers.
Every description and every instance of the schema is finite external
syntax.

The axiom formulation \eqref{eq:eic-induction} and the usual contextual
induction rule are interchangeable.  For example, if
\[
 \Gamma\vdash E(0,\vec p),
 \qquad
 \Gamma,E(z,\vec p)\vdash E(Sz,\vec p),
\]
where $z$ is absent from $\Gamma$ and the passive parameters, discharge the
second premise's equation and universally introduce $z$.  The axiom then
gives $\Gamma\vdash\forall n E(n,\vec p)$, which can be instantiated at any
target term.  Conversely, assume the antecedent of the axiom and apply the
contextual rule with fresh variables.  This uses the ordinary intuitionistic
logical rules and permits arbitrary, including quantified, formulas in
$\Gamma$.

\subsection{Constructor separation and two elementary tests}

Define the sign description and the countdown description by
\begin{equation}
 \operatorname{sg}(u)=C(u,0,1),
 \qquad
 K(u)=C(u,1,Pu).
 \label{eq:eic-sign-countdown}
\end{equation}
Their defining equations imply
\begin{align}
 \operatorname{sg}(0)&=0,
 &\operatorname{sg}(Sx)&=1,\label{eq:eic-sign-computation}\\
 K(0)&=1,
 &K(Sx)&=x.\label{eq:eic-countdown-computation}
\end{align}

\paragraph{Constructor laws.}
Successor is injective: applying $P$ to $Sx=Sy$ gives $x=y$ by
\eqref{eq:eic-predecessor}.  Successor never equals zero.  Indeed,
$Sx=0$ gives $1=0$ by applying $\operatorname{sg}$ and using
\eqref{eq:eic-sign-computation}; symmetry and
\eqref{eq:eic-separation} give $\bot$.  Consequently $1\ne0$, where
$A\ne B$ abbreviates $A=B\to\bot$.  None of these arguments uses
induction.

\paragraph{The sign test.}
For every $u$,
\begin{align}
 &\operatorname{sg}(u)=0\lor\operatorname{sg}(u)=1,
       \label{eq:eic-sign-range}\\
 &\operatorname{sg}(u)=0\quad\longleftrightarrow\quad u=0,
       \label{eq:eic-sign-zero}\\
 &\operatorname{sg}(u)=1\quad\longleftrightarrow\quad u\ne0.
       \label{eq:eic-sign-one}
\end{align}
To prove the range assertion, split \eqref{eq:eic-coverage}.  The branch
$u=0$ gives sign $0$, and the branch $u=S(Pu)$ gives sign $1$.
The implication $u=0\to\operatorname{sg}(u)=0$ is computation.
For its converse, the successor branch of coverage would make
$\operatorname{sg}(u)=0$ imply $1=0$, a contradiction; in that branch
explosion gives $u=0$ as well.  Similarly, if
$\operatorname{sg}(u)=1$ and $u=0$, computation gives $0=1$.
Conversely, $u\ne0$ excludes the zero branch of coverage, and the
successor branch computes the sign as $1$.

Coverage and \eqref{eq:eic-selector} also give, for all $u,v$,
\begin{equation}
 C(u,v,v)=v.
 \label{eq:eic-selector-idempotence}
\end{equation}
Both coverage branches compute to $v$, so this identity again needs no
induction.

\paragraph{The countdown test and the notation $W$.}
For every $u$,
\begin{equation}
 K(u)=0\quad\longleftrightarrow\quad u=1.
 \label{eq:eic-countdown-zero}
\end{equation}
If $u=1=S0$, \eqref{eq:eic-countdown-computation} gives $K(u)=0$.
Conversely, assume $K(u)=0$ and split coverage.  If $u=0$, then
$K(u)=1$, contradicting separation, and explosion gives $u=1$.
If $u=S(Pu)$, computation gives $K(u)=Pu$; hence $Pu=0$ and
$u=S0=1$.

If a presentation uses a description named $W$ with defining equations
\[
 W(0)=1,\qquad W(Sx)=x,
\]
then it proves $K(u)=W(u)$ for every $u$ already from coverage.
In the zero branch both sides are $1$; in the branch $u=S(Pu)$ both
sides are $Pu$.  Thus $K$ and $W$ are provably equal on the whole domain,
even if they are distinct descriptions.  This equality does not appeal to
standard-model agreement or to uniqueness of an arbitrary recurrence.

\subsection{Four elementary uses of equation induction}

Choose descriptions for addition and truncated subtraction, with recursion
on the second argument:
\begin{align}
 x+0&=x, &x+Sy&=S(x+y),\label{eq:eic-addition}\\
 d(x,0)&=x, &d(x,Sy)&=P(d(x,y)).\label{eq:eic-subtraction}
\end{align}
Using these equations and \eqref{eq:eic-induction}, we now derive
\begin{align}
 0+y&=y,\label{eq:eic-left-zero}\\
 (Sx)+y&=S(x+y),\label{eq:eic-left-successor}\\
 d(Sx,Sy)&=d(x,y),\label{eq:eic-cancel-successors}\\
 d(Sa,a)&=1.\label{eq:eic-diagonal}
\end{align}
Each induction below has a single equation as its induction formula.

For \eqref{eq:eic-left-zero}, induct on $y$.  The base is $0+0=0$.
The step is $0+Sy=S(0+y)=Sy$.

For \eqref{eq:eic-left-successor}, hold $x$ fixed and induct on $y$.
At zero both sides are $Sx$.  At a successor the induction hypothesis
gives
\[
 (Sx)+Sy=S\bigl((Sx)+y\bigr)
       =S\bigl(S(x+y)\bigr)=S(x+Sy).
\]
Combining \eqref{eq:eic-left-successor} with the defining equation for
addition gives the useful rearrangement
\begin{equation}
 (Sy)+r=y+Sr.
 \label{eq:eic-shift-successor}
\end{equation}
Indeed, both sides are $S(y+r)$.

For \eqref{eq:eic-cancel-successors}, hold $x$ fixed and induct on $y$.
The base computes as
\[
 d(Sx,S0)=P(d(Sx,0))=P(Sx)=x=d(x,0).
\]
Applying $P$ to the induction equation gives the step:
\[
 d(Sx,S(Sy))=P(d(Sx,Sy))=P(d(x,y))=d(x,Sy).
\]

Finally, induct on $a$ for \eqref{eq:eic-diagonal}.  The base is
$d(S0,0)=S0=1$.  At a successor,
\[
 d(S(Sa),Sa)=d(Sa,a)=1,
\]
using \eqref{eq:eic-cancel-successors} and the induction hypothesis.

\subsection{An equation invariant that reflects actual equality}

The remaining arithmetic ingredient is the following equation, valid for
all $x,y$:
\begin{equation}
 C\bigl(\operatorname{sg}(d(x,y)),\ x,\ y+d(x,y)\bigr)=x.
 \label{eq:eic-residual-invariant}
\end{equation}
Its meaning is that a nonzero subtraction residual reconstructs the
original argument by addition.  The selector expresses this assertion as
one literal equation, so equation induction applies directly.

Hold $x$ fixed and induct on $y$.  At $y=0$, subtraction gives residual
$x$, and \eqref{eq:eic-left-zero} gives $0+x=x$.  The desired equation
is consequently $C(\operatorname{sg}(x),x,x)=x$, which is
\eqref{eq:eic-selector-idempotence}.

For the step, assume \eqref{eq:eic-residual-invariant} at $y$ and
abbreviate
\[
 r=d(x,y),\qquad r'=Pr=d(x,Sy).
\]
These are local term abbreviations, not additional function symbols.
Split coverage for $r'$.

If $r'=0$, then $\operatorname{sg}(r')=0$, and the successor instance
of \eqref{eq:eic-residual-invariant} reduces immediately to $x=x$.

In the other branch, $r'=S(Pr')$, so constructor separation gives
$r'\ne0$.  Split coverage for $r$.  Its zero branch would imply
$r'=P0=0$, a contradiction, from which the desired step equation follows.
In its successor branch,
\[
 r=S(Pr)=Sr',\qquad \operatorname{sg}(r)=1.
\]
The induction equation therefore computes to $y+r=x$.  Meanwhile,
\eqref{eq:eic-shift-successor} gives
\[
 (Sy)+r'=y+Sr'=y+r=x.
\]
Since $r'$ is a successor, $\operatorname{sg}(r')=1$, so this is
exactly the branch selected in the successor instance of
\eqref{eq:eic-residual-invariant}.  This completes the equation induction.
All case distinctions came from coverage; no induction on an implication
or a disjunction was used.

\paragraph{Reflection at residual $1$.}
We have, for arbitrary $a,b$,
\begin{equation}
 d(Sa,b)=1\quad\longleftrightarrow\quad a=b.
 \label{eq:eic-residual-reflection}
\end{equation}
For the forward direction, substitute $x=Sa$, $y=b$ into
\eqref{eq:eic-residual-invariant}.  If $d(Sa,b)=1$, the sign is $1$
and the selector returns its second branch, giving
\[
 b+1=Sa.
\]
The defining equation for addition gives $b+1=S(b+0)=Sb$.  Thus
$Sb=Sa$, and applying $P$ gives $b=a$, hence $a=b$.
Conversely, $a=b$ and \eqref{eq:eic-diagonal} give
$d(Sa,b)=d(Sa,a)=1$.

Define the equality characteristic by the fixed description
\begin{equation}
 e(a,b)=\operatorname{sg}\bigl(K(d(Sa,b))\bigr).
 \label{eq:eic-equality-characteristic}
\end{equation}
The sign, countdown, and residual reflection lemmas give
\begin{align}
 e(a,b)=0&\quad\longleftrightarrow\quad a=b,
       \label{eq:eic-equality-reflection}\\
 e(a,b)=0&\ \lor\ e(a,b)=1.
       \label{eq:eic-equality-range}
\end{align}
In detail, $e(a,b)=0$ is equivalent to $K(d(Sa,b))=0$ by
\eqref{eq:eic-sign-zero}, which is equivalent to $d(Sa,b)=1$ by
\eqref{eq:eic-countdown-zero}, and then to $a=b$ by
\eqref{eq:eic-residual-reflection}.  The range assertion is
\eqref{eq:eic-sign-range} applied to $K(d(Sa,b))$.

The reflection property concerns the logic's actual equality throughout
the domain.  It was not assumed as part of a numerical comparison symbol,
and its proof did not restrict $a$ or $b$ to standard numerals.
It also yields decidable equality.  In the $e(a,b)=0$ branch,
\eqref{eq:eic-equality-reflection} gives $a=b$.  In the $e(a,b)=1$
branch, an assumption $a=b$ would give $e(a,b)=0$ and therefore $0=1$,
contradicting separation.  Hence
\begin{equation}
 \forall a\,\forall b\,\bigl(a=b\lor(a=b\to\bot)\bigr)
 \label{eq:eic-equality-decidable}
\end{equation}
is a theorem, not an extra hypothesis.

\subsection{Compilation of all quantifier-free formulas}

For each fixed quantifier-free formula $A(\vec p)$, construct a term
$b_A(\vec p)$ by external recursion on its finite syntax.  Take equality,
$\bot$, $\land$, $\lor$, and $\to$ as the primitive formula constructors;
negation is implication into $\bot$.  Set
\begin{align}
 b_{s=t}&=e(s,t), & b_{\bot}&=1,\label{eq:eic-boolean-atoms}\\
 b_{A\land D}&=C(b_A,b_D,1),\label{eq:eic-boolean-and}\\
 b_{A\lor D}&=C(b_A,0,b_D),\label{eq:eic-boolean-or}\\
 b_{A\to D}&=C(b_A,b_D,0).\label{eq:eic-boolean-implies}
\end{align}
Every term here is a finite composition of fixed primitive-recursive
descriptions; if desired, projections and composition give it a single
description symbol.  No internal quantification over formulas or programs
is involved.

An external induction on the syntax of $A$ proves, within the arithmetic
theory,
\begin{equation}
 \bigl(b_A=0\lor b_A=1\bigr)
 \quad\text{and}\quad
 \bigl(b_A=0\longleftrightarrow A\bigr).
 \label{eq:eic-boolean-correctness}
\end{equation}
Here are the full connective checks.  In every case $b_A=1$ implies
$\neg A$: if $A$ held, the already established implication
$A\to b_A=0$ would contradict separation.

\paragraph{Atomic formulas and falsity.}
For an equation, both claims are
\eqref{eq:eic-equality-reflection}--\eqref{eq:eic-equality-range}.
For falsity the range assertion is $1=0\lor1=1$, with its right disjunct
proved by reflexivity.  Also $1=0\to\bot$ follows from separation and
symmetry, while $\bot\to1=0$ follows by explosion.  Thus $1=0$ is
equivalent to $\bot$, as required.

\paragraph{Conjunction.}
Split $b_A=0\lor b_A=1$.  In the zero branch, $A$ holds and
$b_{A\land D}=b_D$.  Its range follows from the range of $b_D$,
and its being zero is equivalent to $D$, hence to $A\land D$ under
the known $A$.  In the one branch, $\neg A$ holds and
$b_{A\land D}=1$.  Its range is immediate.  If it were zero,
separation would give $\bot$, hence $A\land D$; conversely an
assumption $A\land D$ contradicts $\neg A$, and explosion gives
$b_{A\land D}=0$.

\paragraph{Disjunction.}
In the branch $b_A=0$, the formula $A$ holds and
$b_{A\lor D}=0$.  The range is immediate, and both $b_{A\lor D}=0$
and $A\lor D$ hold, proving their equivalence.  In the branch
$b_A=1$, we have $\neg A$ and $b_{A\lor D}=b_D$.  Its range follows
from the range of $b_D$.  If its value is zero, then $D$ and hence
$A\lor D$.  Conversely, from $A\lor D$, the $A$ alternative
contradicts $\neg A$ and yields $b_D=0$ by explosion, while the
$D$ alternative yields $b_D=0$ directly.  This proves the required
equivalence in both alternatives.

\paragraph{Implication.}
In the branch $b_A=0$, the formula $A$ holds and
$b_{A\to D}=b_D$.  Its range follows from the range of $b_D$.
If its value is zero, then $D$, which gives $A\to D$.  Conversely,
$A\to D$ and the known $A$ give $D$, hence $b_D=0$.
In the branch $b_A=1$, we have $\neg A$ and
$b_{A\to D}=0$.  Its range is immediate.  To prove $A\to D$,
assume $A$; this contradicts $\neg A$, and explosion gives $D$.
Thus both $b_{A\to D}=0$ and $A\to D$ hold in this branch, and
their equivalence follows.  This argument uses the proved two-value
range of the finite characteristic.  It does not use excluded middle
for arbitrary formulas or change the surrounding intuitionistic logic.

This completes the external syntactic induction proving
\eqref{eq:eic-boolean-correctness} for every quantifier-free formula.

\subsection{Recovery of induction and the full theory}

Let $A(n,\vec p)$ be any fixed quantifier-free formula.  Suppose an
arbitrary retained context $\Gamma$ proves its base and its induction
step:
\[
 \Gamma\vdash A(0,\vec p),
 \qquad
 \Gamma,A(z,\vec p)\vdash A(Sz,\vec p),
\]
with the usual freshness condition on $z$.  By
\eqref{eq:eic-boolean-correctness}, the base gives
$\Gamma\vdash b_A(0,\vec p)=0$.  Under the additional assumption
$b_A(z,\vec p)=0$, the same equivalence gives $A(z,\vec p)$,
the given step gives $A(Sz,\vec p)$, and the equivalence gives
$b_A(Sz,\vec p)=0$.  Apply literal-equation induction to
$b_A(n,\vec p)=0$ and then apply its reflection implication.
The result is $\Gamma\vdash A(t,\vec p)$ for every target term $t$.
Thus the theory proves the full contextual quantifier-free induction
rule, including in arbitrary quantified contexts.

\paragraph{Completion theorem.}
In the common language of fixed primitive-recursive descriptions, the
theory
\[
 \mathsf{PRE}+\mathsf{Cov}+\mathsf{Sep}+\mathsf{EI}
\]
has exactly the same sequents as ordinary first-order intuitionistic
primitive recursive arithmetic with all quantifier-free induction
instances.  The forward simulation follows from the recovered
quantifier-free induction, successor injectivity, and successor separation;
all primitive-recursive defining equations and all logical rules are
shared.  For the reverse simulation, literal-equation induction is an
instance of quantifier-free induction, and successor separation gives
$\mathsf{Sep}$.  Quantifier-free induction proves $\mathsf{Cov}$:
at zero choose its left disjunct, and at a successor use
$Sx=S(P(Sx))$, which follows from $P(Sx)=x$, for its right disjunct.

Both interpretations therefore use the full domain, defined by $x=x$,
and preserve every function symbol, equality, logical connective, and
quantifier.  Both composites have the identity isomorphism, with graph
$x=y$.  The theories are consequently identity-domain bi-interpretable
and equiconsistent.  This completion theorem is unconditional once the
literal-equation induction rule has been obtained; it has no dependency
on a classical successor-uniqueness reduction.


\section{Pairing, primitive recursion, and internal numerical composites}
\label{sec:iterator-coding}
This appendix is adapted with attribution from the progress-closure
arithmetic development~\cite{progress}. Its prerequisites have now been
proved: scalar predecessor, arithmetic reflection, and full quantifier-free
induction. Pairing is a derived proof tool; it is not part of any checker.

\subsection{Pairing from a bounded square-root search}

Once QFI is available, a familiar square-based injection gives a
particularly direct compiler. Write $a\le b$ for $d(a,b)=0$.
The usual linear-order and semiring identities follow by fixed
quantifier-free inductions from the successor, addition, and monus
equations. Multiplication itself is the scalar term
\[
                        a\cdot b=I(b,0;u.u+a).
\]
Only primitive-recursive arithmetic is involved in the following
bounded root search; there is no minimization over an unbounded domain.

For a fixed parameter $n$, put
\[
 U_n(v)=C(d((Sv)^2,n),Sv,v),\qquad
 R_n(i)=I(i,0;v.U_n(v)),\qquad q(n)=R_n(n).
\]
The update increments $v$ exactly when $(v+1)^2\le n$, and
otherwise retains it. QFI in $i$, with $n$ fixed, proves
\[
 R_n(i)\le i,\qquad R_n(i)^2\le n,\qquad
 R_n(i)=i\ \vee\ n<(S(R_n(i)))^2.              \tag{Root invariant}
\]
The base has $R_n(0)=0$. If the update increments, its test
contradicts the strict-upper-bound alternative in the third invariant.
Thus the old state equals $i$, and the new state equals $Si$;
its square bound is the increment test. If the update does not increment, the
strict upper square bound holds and the state remains fixed.
An already established strict upper square bound forces every
subsequent update to remain fixed. These are exactly the cases
needed for the three displayed invariants.

At $i=n$, either the strict upper square bound already holds or
$R_n(n)=n$, in which case $n<(n+1)^2$ supplies it. Thus
\[
                        q(n)^2\le n<(q(n)+1)^2.           \tag{Root}
\]
Monotonicity of squares shows uniqueness of a number satisfying
Root. All formulas just used are quantifier free in the scalar
language, with fixed $n$ during the induction.

Now define
\[
 \pi(a,b)=(a+b)^2+a,\qquad
 L(n)=d(n,q(n)^2),\qquad H(n)=d(q(n),L(n)).
\]
For $k=a+b$, one has $k^2\le k^2+a<(k+1)^2$ since $a\le k$.
Root uniqueness gives $q(\pi(a,b))=a+b$, and therefore
\[
                    L(\pi(a,b))=a,\qquad H(\pi(a,b))=b.
\]
Surjectivity of $\pi$ is unnecessary. The original numerical domain
is never replaced by the range of this coding; pairing is used only
to store the finite state of a fixed primitive recursion. This
standard square pairing is a convenient reconstruction here, not
a claimed new arithmetic basis.

\subsection{Compiling each primitive recursion}

Suppose the fixed target description is
\[
 f(0,\vec p)=g(\vec p),\qquad
 f(Sn,\vec p)=h(n,f(n,\vec p),\vec p).
\]
Assume by external structural induction that the earlier $g,h$
already have scalar translations $g^*,h^*$. Define
\[
 W(u,\vec p)=\pi(SL(u),h^*(L(u),H(u),\vec p)),
\]
\[
 V(n,\vec p)=I(n,\pi(0,g^*(\vec p));u.W(u,\vec p)),
 \qquad f^*(n,\vec p)=H(V(n,\vec p)).
\]
The update has only its state variable bound, with every $p_i$
fixed. Equation induction proves $L(V(n,\vec p))=n$.
The two projection equations and scalar computation then prove
both defining equations for $f^*$ internally. Initial functions
and composition translate directly. This completes the translation
of every fixed finite PR description without adding a universal
evaluator.

\subsection{Why both numerical round trips are internal}

Two fixed scalar iterators with the same seed and provably equal
updates are internally equal at all counters: use literal-equation
induction comparing their values. Their base agrees, and their
successor equations plus congruence prove the step. The same
uniqueness argument works for ordinary primitive recursion.

For an ordinary PR description, induct externally on its finite
construction. The double translation satisfies its original
defining equations, so internal equation induction identifies it
with the original operation. For a native scalar description,
earlier translated updates agree by the external induction
hypothesis; scalar uniqueness identifies the double-translated
iterator. These proofs hold at arbitrary internal counters.

Thus the domain relation $x=y$ is provably a total bijection
preserving every fixed numerical operation in both composites.
The formula constructions and round trips in
Section~\ref{sec:choice-introduction} give the corresponding formula
composites. Both interpretations preserve falsity, taking native
$0=1$ to an ordinary contradiction and ordinary $\bot$ to native
$0=1$. Finite proof translations therefore establish equiconsistency.


% Standalone section for inclusion in a report using amsmath, amssymb,
% amsthm and theorem/lemma/proposition environments.
% New guarded-core argument, developed against ConciseIPRA b17dd5487f34.

\section{A guarded core with a deterministic introduction rule}
\label{sec:guarded-core}

The change is to the public formula language and the introduction rule,
not to the numerical domain.  An atomic matrix immediately below a
paired binder must have a fixed, syntactically recognizable zero guard.
Such a matrix has a provable zero instance.  Consequently it needs only
the progress introduction; the separate atomic case introduction can
be deleted.  Compound matrices retain the case introduction.

\subsection{Complete syntax, axioms, and inference rules}

Numerical terms are those of the inherited scalar calculus:
\[
 t ::= x\mid 0\mid S(t)\mid I(n,a;u.h).
\]
Only the displayed variable $u$ is bound in $h$.  There is one numerical
sort, fixed finite term syntax, and capture-free substitution.  Put
\[
 C(v,a,b):=I(v,a;u.b),\qquad 1:=S0,
\]
where $u$ is fresh.  For a literal equation $E\equiv(s=t)$, define
\[
 G_v E\;:=\;\bigl(C(v,0,s)=C(v,0,t)\bigr).
\tag{Guard}\label{eq:guarded-guard}
\]
The displayed abbreviations expand into terms of the given syntax.

Call a literal equation a \emph{zero-initialized matrix in $y$} when it
has the syntactic form
\[
 B(x,y)\equiv
 I(y,0;u.h)=I(y,0;u.k).
 \tag{Canonical atomic matrix}\label{eq:guarded-canonical-matrix}
\]
Here $h,k$ may depend arbitrarily on the state variable $u$, on $x,y$,
and on fixed parameters; state binders are renamed as necessary.
The selector guard $G_yE$ is a special case, with state-independent
updates.  Every matrix of \eqref{eq:guarded-canonical-matrix} has the
theorem $B(x,0)$ by the two zero computations and equality.

The canonical formulas of $\mathcal G$ are generated as follows:
\begin{enumerate}
\item Every literal equation $s=t$ is a formula.
\item $\mathsf D_{x,y}B$ is a formula when $B$ is already a compound
      canonical formula (its root is $\mathsf D$).
\item $\mathsf D_{x,y}B(x,y)$ is a formula whenever $B$ has the
      zero-initialized shape \eqref{eq:guarded-canonical-matrix}.
\end{enumerate}
The bound variables $x,y$ are distinct.  Rule 3 is the restriction on
atomic matrices.  It does not require an equality decision procedure or
a proof that an arbitrary base is true: the guard has the displayed
literal shape.  The designated contradiction is $0=1$.

Retained contexts contain reusable clauses
$\langle A;y.B\rangle$, where $A,B$ are canonical formulas and
$y\notin\operatorname{FV}(A)$.  An ordinary assumption $F$ abbreviates
the constant clause $\langle H;v.F\rangle$, where $v$ is fresh and
$H$ is the fixed computation theorem $I(0,0;u.0)=0$.

The two nonlogical axiom schemas, available in every context, are
\[
 I(0,a;u.h)=a,\qquad
 I(Sn,a;u.h)=h[I(n,a;u.h)/u].
\tag{I0,IS}\label{eq:guarded-computation}
\]
Besides these, there are precisely the following rules:
\[
 \frac{\Gamma\vdash A}{\Gamma\vdash B(t)}
 \quad\text{if }\langle A;y.B(y)\rangle\in\Gamma
 \qquad(\mathrm{Use}),
\]
\[
 \frac{\Gamma\vdash s=t\qquad\Gamma\vdash F(s)}
      {\Gamma\vdash F(t)}\qquad(\mathrm{Tr}),
\]
\[
 \frac{\Gamma\vdash\mathsf D_{x,y}B(x,y)\qquad
       \Gamma,\langle B(a,0);y.B(a,y)\rangle\vdash F}
      {\Gamma\vdash F}\qquad(\mathrm{DE}),
\]
and the single, root-determined introduction schema
\[
 \frac{\Gamma,B(t,\eta_B(z))\vdash B(t,Sz)}
      {\Gamma\vdash\mathsf D_{x,y}B(x,y)},
 \qquad
 \eta_B(z)=
 \begin{cases}
 z,& B\text{ is a literal equation},\\
 0,& B\text{ is compound}.
 \end{cases}
 \tag{DI$_{\rm root}$}\label{eq:guarded-di}
\]
There is no certificate choice between these two cases.  At an atomic
matrix the guard condition on the conclusion is mandatory in the
canonical calculus.  In DE, $a$ is fresh for $\Gamma,F$ and the fixed
parameters.  In DI, $z$ is fresh for $\Gamma,t$ and the fixed
parameters.  The displayed local resources are discharged.
Transport ranges over canonical formula templates, and capture-free
numerical substitution preserves their guard shape.

There are no additional primitive logical rules, reflexivity axioms,
successor axioms, coverage axioms, or induction schemas.
Weakening, exchange, contraction, renaming, numerical substitution, and
ordinary-assumption cut are admissible by induction on derivations.
For cut, replace a use of the corresponding constant clause by the
supplied proof, renaming eigenparameters first.

\subsection{Padding, equation implication, and numerical cases}

Write $A\dashv\vdash B$ for interderivability under arbitrary additional
retained context.  The following preliminary arguments use no
contextual induction rule.

\begin{lemma}[Equality and guard endpoints]
Reflexivity, symmetry, transitivity, numerical congruence, and their
usual contextual uses are derivable.  Moreover,
\[
 \vdash G_0 E,\qquad G_{Sz}E\dashv\vdash E.
\]
\end{lemma}
\begin{proof}
For reflexivity, put $r=I(0,t;u.0)$.  The computation equation $r=t$
serves as both premises of transport with motive $v=t$, yielding
$t=t$.  Standard equality motives then give the remaining equality
rules.  The two guard claims follow from the two computation equations
for $C$, together with equality transport.
\end{proof}

For an equation $E$ let
\[
 L(E):=\mathsf D_{u,v}G_v E,
 \qquad
 P(A):=
 \begin{cases}L(A),&A\text{ is an equation},\\A,&A\text{ is compound}.
 \end{cases}
\tag{Padding}\label{eq:guarded-padding}
\]
The newly bound $u,v$ are fresh for $E$.  In particular, $P(A)$ is
always compound.

\begin{lemma}[Padding]
$L(E)\dashv\vdash E$, and therefore $P(A)\dashv\vdash A$.
\end{lemma}
\begin{proof}
Under $E$, atomic progress introduction derives $L(E)$: its target
$G_{Sz}E$ follows from the retained $E$, independently of the local
induction assumption.  Conversely, DE on $L(E)$ opens
$\langle G_0E;v.G_vE\rangle$.  Supply the guard-zero theorem and use
the clause at $1$; guard removal gives $E$.
\end{proof}

We will need to reason by numerical cases under an assumption that
contains the case variable.  It is essential not to violate the
eigenvariable condition by directly retaining that assumption during
case introduction.  A small implication macro resolves this issue.

For equations $U\equiv(s=t)$ and $V\equiv(r=q)$ put
\[
 M_w(U,V):=(C(w,s,r)=C(w,t,q)),\qquad
 U\Rightarrow_g V:=\mathsf D_{a,w}L(M_w(U,V)),
\tag{Equation implication}\label{eq:guarded-implication}
\]
with fresh padding variables and unused $a$.
The matrix of the outer binder is compound.

\begin{lemma}[Implication between equations]
The macros in \eqref{eq:guarded-implication} admit both rules
\[
 \frac{\Gamma,U\vdash V}{\Gamma\vdash U\Rightarrow_g V},
 \qquad
 \frac{\Gamma\vdash U\Rightarrow_g V\qquad\Gamma\vdash U}
      {\Gamma\vdash V}.
\]
\end{lemma}
\begin{proof}
The endpoint identities for $C$ and padding give
$L(M_0(U,V))\dashv\vdash U$ and
$L(M_{Sz}(U,V))\dashv\vdash V$.
Compound case DI therefore gives introduction.  DE, the first endpoint,
Use at $1$, and the second endpoint give elimination.  These endpoint
conversions use numerical transport with parameters outside the fresh
padding binders; they do not assume a binder-congruence theorem.
\end{proof}

\begin{lemma}[Numerical cases]
For every canonical $A(y)$ the following rule is admissible:
\[
 \frac{\Gamma\vdash A(0)\qquad\Gamma\vdash A(Sz)}
      {\Gamma\vdash A(t)},
 \qquad z\text{ fresh for }\Gamma\text{ and the fixed parameters}.
 \tag{NC}\label{eq:guarded-cases}
\]
\end{lemma}
\begin{proof}
First suppose $A$ is compound.  Case DI with an unused witness coordinate
introduces $\mathsf D_{u,y}A(y)$, ignoring its local base assumption in
the successor premise.  DE opens $\langle A(0);y.A(y)\rangle$; supply
the given base proof and use it at $t$.

If $A$ is an equation, apply this already proved compound case rule to
$L(A(y))$.  Padding converts the two premises to its endpoints and
converts the conclusion back.  No progress induction is being used in
this construction beyond the special padding introductions.
\end{proof}

\subsection{Recovery of contextual equation induction}

\begin{theorem}[Equation induction]
For every literal equation $E(y)$, $\mathcal G$ derives
\[
 \frac{\Gamma\vdash E(0)\qquad\Gamma,E(z)\vdash E(Sz)}
      {\Gamma\vdash E(t)},
 \qquad z\text{ fresh for }\Gamma\text{ and the fixed parameters}.
 \tag{EI}\label{eq:guarded-ei}
\]
\end{theorem}
\begin{proof}
Put $J(y)=G_yE(y)$.  We first prove the progress premise
$\Gamma,J(z)\vdash J(Sz)$ without making an invalid contextual use of
NC.  Consider the compound formula
\[
 \Phi(y):=J(y)\Rightarrow_g E(Sy).
\]
Under $\Gamma$, prove $\Phi(0)$ by equation-implication introduction:
ignore its local hypothesis, use the supplied proof of $E(0)$, and
apply the supplied step proof specialized at $0$ to obtain $E(1)$.

For fresh $w$, prove $\Phi(Sw)$ by equation-implication introduction.
Its local hypothesis $J(Sw)$ gives $E(Sw)$ by guard removal, and the
supplied step proof specialized at $Sw$ gives $E(SSw)$.
These two branch proofs have retained context $\Gamma$, independent
of the case variable.  Thus NC legitimately gives $\Phi(z)$ under
$\Gamma$.  Under the local hypothesis $J(z)$, equation-implication
elimination gives $E(Sz)$, hence $J(Sz)$ by the successor guard endpoint.

Atomic progress DI now introduces $\mathsf D_{u,y}J(y)$ with unused
$u$.  DE and the theorem $J(0)$ yield $J(Sw)$, and hence $E(Sw)$, for
a fresh $w$.  Finally NC for $E$ uses the supplied base $E(0)$ and
these successor instances to conclude $E(t)$.

The only induction-like premise used here is the single displayed
guarded progress introduction.  In particular, no predecessor,
coverage disjunction, formula induction, or pre-existing EI is assumed.
\end{proof}

\subsection{Exact translations to and from the inherited calculus}

Let $\mathcal C$ be the inherited Choice-DI calculus: the same scalar
terms, arbitrary formulas $s=t$ and $\mathsf D_{x,y}B$, the same
computation, Use, DE, and transport rules, with case DI at every matrix
and progress DI at literal equations.  Define $T:\mathcal C\to\mathcal G$
on finite formulas by
\[
 T(s=t)=(s=t),\qquad
 T(\mathsf D_{x,y}B)=\mathsf D_{x,y}P(TB).
\tag{Translation}\label{eq:guarded-translation}
\]
Each padding binder uses fresh names.  This is a computable structural
translation, preserves free numerical variables, commutes with
capture-free numerical substitution up to bound-variable renaming,
and translates every paired matrix to a compound matrix.
Translate clauses by
$T\langle A;y.B\rangle=\langle TA;y.TB\rangle$.
Let $J:\mathcal G\to\mathcal C$ be literal inclusion.

\begin{theorem}[Translation of derivations]
\[
 \Gamma\vdash_{\mathcal C}A\Longrightarrow
 T\Gamma\vdash_{\mathcal G}TA,
 \qquad
 \Gamma\vdash_{\mathcal G}A\Longrightarrow
 J\Gamma\vdash_{\mathcal C}JA.
\]
The transformations are effective on finite derivations.
\end{theorem}
\begin{proof}
Inclusion is immediate: every guarded rule is an allowed Choice-DI
rule.  For $T$, induct on a Choice-DI derivation.

Computation is unchanged.  Use commutes with substitution, and
transport uses the translated motive.  A case DI premise translates
to a proof from $TB(t,0)$ to $TB(t,Sz)$.  Under the local assumption
$P(TB)(t,0)$, unpad, apply that translated proof, and pad the result.
Compound case DI then proves the required
$\mathsf D_{x,y}P(TB)$.

For a progress DI with equation matrix $E$, its translated step is
still $\Gamma,E(t,z)\vdash E(t,Sz)$.  To introduce
$\mathsf D_{x,y}L(E(x,y))$, use compound case DI.  Its retained base
$L(E(t,0))$ supplies $E(t,0)$.  The recovered contextual EI, under
that base, gives $E(t,Sz)$, which padding converts to the target
$L(E(t,Sz))$.

For DE, the translated major is $\mathsf D_{x,y}P(TB)$.
It opens the clause
$\langle P(TB)(a,0);y.P(TB)(a,y)\rangle$.
The recursively translated minor expects
$\langle TB(a,0);y.TB(a,y)\rangle$.  Replace every use of this
expected clause by: pad its antecedent, use the opened clause at the
same term, and unpad its consequent.  This effective local proof
replacement preserves the fresh witness condition and the conclusion.
\end{proof}

\begin{lemma}[Compound-matrix congruence]
If canonical compound matrices $B,C$ are uniformly interderivable,
then $\mathsf D_{x,y}B\dashv\vdash\mathsf D_{x,y}C$ in $\mathcal G$.
\end{lemma}
\begin{proof}
Open the source binder at a fresh witness $a$.  Introduce the target
by compound case DI at that witness.  Convert its retained base to
the source base, use the opened source clause at $Sz$, and convert
the successor instance to the target instance.  Close DE.  The
reverse conversion proves the converse.
\end{proof}

\begin{lemma}[The guarded atomic composite]
For an admissible zero-initialized atomic matrix $B(x,y)$,
\[
 \mathsf D_{x,y}B\dashv\vdash\mathsf D_{x,y}L(B)
 \quad\text{in }\mathcal G.
 \tag{Atomic composite}\label{eq:guarded-atomic-composite}
\]
\end{lemma}
\begin{proof}
Forwards, open the source clause at $a$.  Introduce the target by
compound case DI: unpad its base, use the source clause at $Sz$, and
pad the result.

Backwards, open the clause of $\mathsf D_{x,y}L(B)$ at $a$.
The two zero computations prove $B(a,0)$, and padding proves
$L(B(a,0))$.
Thus the opened clause gives $L(B(a,Sz))$, and unpadding gives
$B(a,Sz)$, independently of the local assumption $B(a,z)$.
Atomic progress DI therefore introduces $\mathsf D_{x,y}B$ at
witness $a$.  Close DE.  This argument uses the true base explicitly;
it would fail for an arbitrary unguarded atomic matrix.
\end{proof}

\begin{theorem}[Both composites and bi-interpretability]
For all $\mathcal C$-formulas $A$ and all canonical
$\mathcal G$-formulas $B$,
\[
 A\dashv\vdash_{\mathcal C}JT(A),\qquad
 B\dashv\vdash_{\mathcal G}TJ(B).
\]
Consequently $\mathcal C$ and $\mathcal G$ admit effective structural
intertranslations preserving their full consequence relations, with
both native formula composites and literal identity numerical
translations.
\end{theorem}
\begin{proof}
In $\mathcal C$, padding is interderivable with its argument, and
case introduction proves binder congruence for arbitrary matrices.
Structural induction gives the first displayed equivalence.

For the second, induct on a canonical $B$.  Equations are unchanged.
At a compound matrix, the induction hypothesis and compound-matrix
congruence apply.  At a guarded atomic matrix, the statement is
exactly \eqref{eq:guarded-atomic-composite}.  The equivalences are
uniform under retained assumptions, so they extend to clauses and
sequent contexts by the same local replacement of Use instances.

Both interpretations leave the numerical sort, every term, and
numerical equality untouched.  Both numerical composites therefore
have the identity graph $x=y$; no quotient, retract, numeral restriction,
or extra coding sort is involved.
\end{proof}

For the first-order bi-interpretability statement, use the ordinary
realization that reads every canonical
$\mathsf D_{x,y}B$ uniformly as
$\exists x(B(x,0)\to\forall yB(x,y))$ and leaves numerical equality
and all scalar terms unchanged.  Composing the displayed translations
and composites with the inherited full-domain bridge between
$\mathcal C$ and iPRA yields a bi-interpretation of this ordinary
realization of $\mathcal G$ with full first-order iPRA.
The contradiction $0=1$ is
unchanged, so equiconsistency follows as well.  The inherited arithmetic
bridge supplies the primitive-recursive compiler and the ordinary
logical interpretation; this section does not claim those inherited
constructions as new.

\subsection{Why an unrestricted checker can still be sound}

There is a useful, but strictly separate, administrative calculus
$\mathcal U$.  It permits arbitrary atomic paired matrices and uses
the same root-determined DI, with ordinary Q-style clauses on DE.
It is not claimed that all of $\mathcal U$ is bi-interpretable with iPRA.

Give arbitrary administrative formulas their following ordinary
interpretations:
\[
 (s=t)^\dagger=(s=t),
\]
\[
 (\mathsf D_{x,y}B)^\dagger=
 \begin{cases}
 \exists x\,\forall y\,(B(x,y)\to B(x,Sy)),
       &B\text{ is a literal equation},\\[2pt]
 \exists x\,(B^\dagger(x,0)\to\forall y\,B^\dagger(x,y)),
       &B\text{ is compound}.
 \end{cases}
 \tag{Mixed administrative reading}\label{eq:guarded-mixed}
\]
Retained clauses keep the ordinary reading
$A^\dagger\to\forall y\,B^\dagger$.
If the raw codec also permits administrative $Q$-closures as
formula-shaped objects, including under transport or a paired binder,
extend this interpretation by
$(Q_yB)^\dagger=B^\dagger(0)\to\forall yB^\dagger(y)$.
Such a matrix is compound for root-determined DI.  This adds no
canonical object-formula constructor and preserves the following
soundness argument.

Every $\mathcal U$ inference is sound in iPRA under this reading.
Atomic DI directly proves its progressive reading; compound case DI
uses numerical coverage.  In atomic DE, contextual equation induction
turns $\forall y(E(y)\to E(Sy))$ into
$E(0)\to\forall yE(y)$, exactly the clause that DE supplies.
All other rules are immediate, including transport: numerical
substitution does not change whether a matrix is an equation.

On canonical zero-initialized matrices, the mixed and ordinary readings
agree provably in iPRA.  Indeed $B(x,0)$ is a theorem, so, with $x$ fixed,
\[
 \forall y\,(B(x,y)\to B(x,Sy))
 \quad\longleftrightarrow\quad
 B(x,0)\to\forall y\,B(x,y).
\]
The forward direction is equation induction; the reverse uses the
true base to obtain every instance.  Structural induction extends
this agreement to all canonical guarded formulas.

Thus administrative proofs with unguarded intermediate formulas are
conservative on canonical guarded conclusions: translate an accepted
derivation into iPRA, use the canonical agreement, and translate back
through the established bi-interpretation.  Conversely, every canonical
guarded derivation is already an administrative derivation.

This justifies a checker contract consisting of soundness and
completeness for canonical guarded statements while allowing broader
administrative certificates.  It does \emph{not} make the raw checker
a validator of the guarded public grammar.  It also does not assert
completeness for every noncanonical encoding under the mixed reading.
If a comparison requires rejection of all noncanonical public goals,
the corresponding syntax validator must be implemented and counted.

\subsection{Why merely deleting the choice would not suffice}

Without the guarded canonical language, root-deterministic DI does
not preserve the original native consequence relation.  Let
\[
 p(y)=I(y,0;u.C(u,1,0)),\qquad E(y)\equiv(p(y)=1),
\]
and leave the existential witness coordinate unused.  Computation
gives $p(0)=0,p(1)=1,p(2)=0$.  The original case introduction proves
$\mathsf D_{x,y}E(y)$: its base gives $0=1$, from which equality
transport and selector computation prove every equation.

By contrast, the progressive reading of this formula is false in the
standard natural numbers, because $E(1)$ holds and $E(2)$ does not.
The mixed interpretation proves that it is not a theorem of the
unrestricted root-deterministic calculus.  The guard, the explicit
translations, and both composite proofs are therefore substantive
parts of the reduction, rather than optional notation.


% Standalone mathematical module; requires amsmath, amssymb, amsthm, url.
% The surrounding document should supply theorem, lemma, and corollary environments.
% Written derivation, not a proof-assistant verification.
\section{Removing literal zero from the source of transport}
\label{zs:section}

The following is a restriction of a mathematical inference rule, rather
than a change in the spelling of zero. Throughout this section,
``nonzero source'' means \emph{not syntactically the single term $0$}.
Every variable and every term headed by $S$ or $I$ is a permitted source,
including terms whose numerical value is zero.

Let $\mathcal C$ be the native scalar calculus with terms
\[
 t::=x\mid0\mid St\mid I(n,a;u.h),
\]
the two computation schemas
\[
 I(0,a;u.h)=a,\qquad
 I(Sn,a;u.h)=h[I(n,a;u.h)/u],
\]
formulas generated by equality and the $Q,D$ binders, and the Hyp, QE,
DE, case-DI and equation-progress-DI rules stated in the main inventory.
The same argument applies to the original presentation with $D$ object
formulas and $Q$ resources in contexts. Let $\mathcal C^{-}$ replace
unrestricted transport by
\[
 \frac{\Gamma\vdash s=t\qquad\Gamma\vdash F(s)}
      {\Gamma\vdash F(t)}
 \quad\mathsf{Tr}^{-},\qquad s\not\equiv0.
 \tag{Tr$^-$}\label{zs:restricted-rule}
\]
Here $F$ is an arbitrary native numerical template, and $s\not\equiv0$
is a condition on finite term syntax. All other rules are unchanged.

\begin{theorem}[Literal-zero source elimination]
\label{zs:main}
For every native context $\Gamma$ and formula $A$,
\[
 \Gamma\vdash_{\mathcal C^{-}}A
 \quad\Longleftrightarrow\quad
 \Gamma\vdash_{\mathcal C}A.
\]
In particular, the ordinary whole-domain realization of
$\mathcal C^{-}$ is bi-interpretable and equiconsistent with full
first-order iPRA through the same interpretations as the native scalar
core. No change of numerical domain or interpreted equality is needed.
\end{theorem}

We prove the theorem without assuming arbitrary numerical substitution
admissibility for $\mathcal C^{-}$: such substitution can turn a variable
source into the forbidden literal zero. Every occurrence of a helper
equation at an arbitrary term is either obtained directly from its
displayed derived rule at that term, or from a derivation whose transport
sources retain an $S$ or $I$ head under the indicated substitution.

\subsection{Structural operations, cases, and equation induction}

Weakening, ordinary-assumption cut, and capture-free renaming are
admissible by induction on derivations. These operations do not turn a
nonzero source into literal zero. As usual, fresh parameters are renamed
before a cut. We also use the following restricted substitution fact:
a capture-free substitution is admissible for a particular finite
derivation if it leaves every primitive transport source in that
derivation syntactically different from zero. This fact follows by the
same induction; it is not unrestricted substitution admissibility.

\begin{lemma}[Cases and equation induction need no transport]
\label{zs:cases-induction}
In $\mathcal C^{-}$ the following rules are derivable:
\[
 \frac{\Gamma\vdash A(0)\qquad\Gamma\vdash A(Sz)}
      {\Gamma\vdash A(t)}\quad\mathsf{NC},
 \qquad
 \frac{\Gamma\vdash E(0)\qquad\Gamma,E(z)\vdash E(Sz)}
      {\Gamma\vdash E(t)}\quad\mathsf{EI}.
\]
In NC, $A$ is any native formula; in EI, $E$ is one literal equation.
The parameter $z$ is fresh for $\Gamma$, $t$, and the fixed parameters
of the matrix. Other parameters stay fixed throughout EI.
\end{lemma}
\begin{proof}
Choose an unused witness coordinate $x$. In NC, weaken the successor
proof by $A(0)$ and apply case-DI, with witness zero, to obtain
$D_{x,y}A(y)$. DE opens $Q_yA(y)$, because $x$ is unused. QE with the
given base proof and target $t$ yields $A(t)$, and DE closes.
For EI use progress-DI instead, with the given equation step, and use
the same DE/QE construction. No transport rule occurs.
\end{proof}

These rules will be applied to unconditional helper equations before
they are weakened into the context containing $0=t$. We never use NC
with a case parameter occurring in its retained context.

\subsection{The equality operations available immediately}

\begin{lemma}[Reflexivity and safe equality operations]
\label{zs:safe-equality}
Reflexivity $r=r$ is derivable for every term $r$. From $p=q$ with
$p\not\equiv0$, one can derive $q=p$ and $H(p)=H(q)$ for any numerical
template $H$. Every direct computation can be used in either direction
as an equation; using the reversed equation as a transport major is
permitted only when its new source is syntactically nonzero.
\end{lemma}
\begin{proof}
Put $b=I(0,r;u.0)$. The computation $b=r$, used as both major and
minor of transport with template $v=r$, gives $r=r$. Its source is
the nonempty iterator term $b$.
For symmetry of $p=q$, transport the reflexivity $p=p$ through the
template $v=p$. For numerical congruence, transport the reflexivity
$H(p)=H(p)$ through $H(p)=H(v)$. Both use the permitted source $p$.
Every computation equation has an $I$-headed left side, so its
symmetry follows by the same construction.
\end{proof}

An occurrence-selective template may keep some occurrences of a term
fixed while abstracting other occurrences. This is ordinary
capture-free template substitution, not an extra inference rule.

\begin{lemma}[Joining two nonzero terms known to be zero]
\label{zs:zero-join}
Suppose $p,q\not\equiv0$ and $\Gamma$ proves $p=0$ and $q=0$.
Then $\Gamma\vdash p=q$ using only $\mathsf{Tr}^{-}$ and computation.
\end{lemma}
\begin{proof}
Let
\[
 z_0=I(0,0;u.0),\qquad b_p=I(0,p;u.0).
\]
Computation gives $b_p=p$. Start with $b_p=b_p$ and use the major
$p=0$ in the template $I(0,v;u.0)=b_p$. This gives $z_0=b_p$.
Symmetry at the nonzero source $z_0$ gives $b_p=z_0$.
Now transport through $v=z_0$ using the computation $b_p=p$, whose
source is $b_p$, to obtain $p=z_0$.
The identical construction gives $q=z_0$. Symmetry at source $q$
gives $z_0=q$, and transport at source $z_0$ in $p=z_0$ gives $p=q$.
The sources used are $p,q,z_0,b_p,b_q$; all are syntactically nonzero.
No reverse transport from zero has been used.
\end{proof}

The zero-join lemma is deliberately placed before the next construction.
The positive case of its singleton test needs precisely this lemma;
ordinary transitivity through zero is not yet available.

\subsection{A singleton-one test without predecessor}

Define the finite scalar terms
\[
\begin{aligned}
 C(v,a,b)&:=I(v,a;u.b),& Z(v)&:=C(v,1,0),\\
 K(v)&:=C(v,0,1),& H(v)&:=C(v,1,2),\\
 N(n)&:=I(n,0;u.Z(u)),& R(n)&:=I(n,0;u.H(u)),\\
 L(n)&:=N(R(n)).
\end{aligned}
\tag{Tests}\label{zs:tests}
\]
The update variable in $C$ is unused. Here $1=S0$ and $2=S1$.
The standard-model description of $R$ is $0,1,2,2,\ldots$; the proof
below establishes the needed equations internally and does not assume
this description.

\begin{lemma}[Uniform test equations]
\label{zs:singleton}
The following equations have derivations in $\mathcal C^{-}$,
uniformly in the displayed argument:
\[
 H(v)=S K(v),\qquad R(1)=1,\qquad R(SSv)=2,
\]
\[
 N(1)=1,\quad N(2)=0,\quad
 L(1)=1,\quad L(SSv)=0,\quad
 \boxed{L(Sv)=Z(v)}.
 \tag{Singleton}\label{zs:singleton-equations}
\]
All transport sources in these helper derivations have an $I$ or $S$
head, except for the explicit nonzero endpoints in the zero-join
lemma; those endpoints also have an $I$ head in this application.
\end{lemma}
\begin{proof}
For $H(v)=S K(v)$ use NC. At zero both sides reduce to $1$; at a
successor both reduce to $2$. To join two equations $p=j$, $q=j$
with $j=1$ or $j=2$, reverse $q=j$ at its nonzero source $q$, then
transport at the nonzero numeral $j$. Each reduction of these
particular terms is a forward replacement of an iterator computation.
This proves the helper at an arbitrary target $v$ by NC, without
substituting zero into a variable-source transport.

The equations $N(1)=1$ and $N(2)=0$ are finite computations:
$N(0)=0$, $N(1)=Z(N(0))=1$, and $N(2)=Z(N(1))=0$.
Every replacement source in these calculations is an iterator term.
Likewise $R(1)=H(R(0))=H(0)=1$.

Computation and the first helper give, for any $v$,
\[
 R(Sv)=H(R(v))=S K(R(v)).
\]
Use this equation at the nonzero source $R(Sv)$ in
$R(SSv)=H(R(Sv))$. The resulting right side is
$H(S K(R(v)))=2$ by the successor computation of $C$.
Thus $R(SSv)=2$. These substitutions are all at $I$-headed sources.
Substitute $R(1)=1$ and $R(SSv)=2$ inside $N$ and use the two
computed values of $N$ to obtain $L(1)=1$ and $L(SSv)=0$.

Finally apply NC to the matrix $L(Sv)=Z(v)$.
At zero, $L(1)=1$ and $Z(0)=1$ are joined through the nonzero numeral
$1$. At a successor, $L(SSv)=0$ and $Z(Sv)=0$ are joined by
Lemma~\ref{zs:zero-join}; both left sides are nonempty iterator terms.
This proves the boxed equation at every target. Its use is therefore
an instance of a displayed derived rule, not an appeal to unrestricted
substitution admissibility.
\end{proof}

\begin{lemma}[An unconditional zero-filter identity]
\label{zs:filter}
For every term $t$,
\[
                 t=C(Z(t),t,0)
                 \tag{Filter}\label{zs:filter-equation}
\]
is derivable in the empty context of $\mathcal C^{-}$.
\end{lemma}
\begin{proof}
Apply NC to $v=C(Z(v),v,0)$, with the target $t$.
At zero, forward computations give
$C(Z(0),0,0)=C(1,0,0)=0$; reverse this equation at its nonzero
left side. At $Sv$, forward computations give
$C(Z(Sv),Sv,0)=C(0,Sv,0)=Sv$; again reverse at the nonzero
iterator left side. Thus the base and successor proofs have no
retained assumption depending on the case variable. All numerical
arguments of the final NC target are permitted. In particular, this
identity is proved before any assumption $0=t$ is used.
\end{proof}

\subsection{Turning a zero-source equality around}

\begin{lemma}[Zero-source symmetry]
\label{zs:zero-symmetry}
For every term $t$,
\[
                     0=t\ \vdash_{\mathcal C^{-}}\ t=0.
                     \tag{ZS}\label{zs:zero-symmetry-equation}
\]
\end{lemma}
\begin{proof}
If $t\equiv0$, use reflexivity. Otherwise $t$ is a permitted source.
Work in the fixed context $\Gamma=\{0=t\}$ and put
\[
                  A(n):=I(n,t;u.Su).
\]
We first prove $n=A(n)$ by EI, keeping $t$ fixed.

For the base, computation gives $A(0)=t$. Its symmetry gives
$t=A(0)$. Since $t\not\equiv0$, transport at source $t$ in the
retained assumption $0=t$ gives $0=A(0)$.

For the step, retain $z=A(z)$ with $z$ fresh. This source is a
variable, hence is permitted. Numerical successor congruence from
Lemma~\ref{zs:safe-equality} gives $Sz=S A(z)$.
Computation gives $A(Sz)=S A(z)$, and its symmetry gives
$S A(z)=A(Sz)$. Its source $S A(z)$ is also permitted, so transport
in $Sz=S A(z)$ gives $Sz=A(Sz)$. This is exactly the EI step.
The induction hypothesis is used at the same fixed parameter $t$.

At target $1$, EI gives $1=A(1)$. Forward computation gives
$A(1)=St$: the successor equation produces $S A(0)$, and the
nonzero-source computation $A(0)=t$ is used inside $S$.
Transport at source $A(1)$ consequently gives
\[
                              1=St.
\]
Its source $1$ is permitted. Numerical congruence for $L$ gives
$L(1)=L(St)$. Use $L(1)=1$ at source $L(1)$, and
$L(St)=Z(t)$ at source $L(St)$, obtaining $1=Z(t)$.
Symmetry at source $1$ gives $Z(t)=1$.

Weaken the unconditional Filter identity into $\Gamma$.
Transport at the iterator source $Z(t)$ changes its right side to
$C(1,t,0)$. The latter computes to zero, and its source is again
an iterator term. The resulting equation is $t=0$, as required.
\end{proof}

This proof uses cases only for unconditional helper equations. The
sole induction with the retained assumption $0=t$ has the literal
equation $n=A(n)$ as invariant. It does not perform a forbidden
case distinction while silently changing the retained context.

\subsection{Recovering equality algebra and formula transport}

\begin{lemma}[Full numerical equality algebra]
\label{zs:full-equality}
The calculus $\mathcal C^{-}$ admits unrestricted symmetry,
numerical congruence, and transitivity. In particular it admits
root replacement
\[
 \frac{\Gamma\vdash s=t\qquad\Gamma\vdash s=r}
      {\Gamma\vdash t=r}.
\]
\end{lemma}
\begin{proof}
Symmetry uses Lemma~\ref{zs:safe-equality} when the source is not
zero and Lemma~\ref{zs:zero-symmetry} otherwise.

For numerical congruence from $s=t$, a nonzero $s$ is handled by
Lemma~\ref{zs:safe-equality}. If $s\equiv t\equiv0$, use
reflexivity. If $s\equiv0$ and $t\not\equiv0$, first derive $t=0$,
apply permitted numerical congruence at source $t$ to obtain
$H(t)=H(0)$, and use the now established full symmetry.

For transitivity from $r=s$ and $s=t$, a nonzero middle term $s$
is handled by one permitted transport. If $s\equiv0$ and either
endpoint is literally zero, a given premise is already the
conclusion. Otherwise $r,t\not\equiv0$; reverse $0=t$ to get
$t=0$ and apply the zero-join lemma to $r=0$ and $t=0$.
Root replacement follows by symmetry and transitivity.
These constructions are acyclic: zero-source symmetry was proved
before any unrestricted numerical congruence or transitivity.
\end{proof}

\begin{lemma}[Full transport]
\label{zs:full-transport}
For every native numerical template $F(v)$, the rule
\[
 \frac{\Gamma\vdash s=t\qquad\Gamma\vdash F(s)}
      {\Gamma\vdash F(t)}
\]
is admissible in $\mathcal C^{-}$, without a source restriction.
\end{lemma}
\begin{proof}
Use external induction on the finite number of formula binders in
$F$, uniformly in retained contexts and fixed numerical parameters.
For an equation $F(v)\equiv A(v)=B(v)$, use unrestricted numerical
congruence to obtain $A(s)=A(t)$ and $B(s)=B(t)$. Join these with
$A(s)=B(s)$ using the established equality algebra.

For $F(v)=D_{x,y}B(v,x,y)$, eliminate $F(s)$ at a fresh $a$,
retaining $Q_yB(s,a,y)$. Introduce the desired $F(t)$ by case-DI
with witness $a$. Its local assumption is $B(t,a,0)$.
Full symmetry gives $t=s$, and the external induction hypothesis
for the smaller matrix transports the assumption to $B(s,a,0)$.
QE gives $B(s,a,Sz)$, and the induction hypothesis in the forward
direction gives $B(t,a,Sz)$. Close case-DI and DE. All new
parameters can be chosen fresh for the source equality and the
retained context.

If $Q$ is first-class formula syntax, first observe that its case
introduction is derived without transport: introduce
$D_{u,y}B(y)$ with unused $u$ by case-DI and eliminate it with the
one-node continuation $Q_yB(y)$. For $F(v)=Q_yB(v,y)$, use this
derived introduction for the target. Convert its base assumption
backwards by the induction hypothesis, apply the old Q resource,
and convert the successor conclusion forwards. This is the same
smaller-matrix argument and uses no new numerical induction.
\end{proof}

\begin{proof}[Proof of Theorem~\ref{zs:main}]
Every restricted derivation is an unrestricted one. Conversely,
replace each unrestricted transport node by the derivation supplied
by Lemma~\ref{zs:full-transport}, recursively translating its
premises first. Other inference nodes are unchanged. The result has
the same context and conclusion. Thus the full native consequence
relation, not just closed theorems, is preserved. The existing
native/ordinary and ordinary/iPRA interpretations and both identity
composites therefore transfer unchanged.
\end{proof}

\subsection{Compatibility with restricted computation certificates}
\label{zs:certificate-computation}

Some compact checkers require the initializer of a computation packet
to have nonempty syntax, and encode a successor counter by a nonempty
predecessor. These are certificate restrictions, not restrictions
on the two mathematical computation schemas. They do not introduce
a circular dependence on zero-source transport in the preceding proof.

Use the inherited certified-zero term
\[
                         c_0=I(2,1;u.0).
\]
Its equation $c_0=0$ is a directly admitted successor computation:
the predecessor and initializer are both $1$. For an arbitrary
computation instance with literal-zero initializer or literal-zero
successor predecessor, form its whole equation as a template $J(v)$,
replacing just these offending occurrences by $v$. The instance
$J(c_0)$ is a directly admitted computation. Transport through this
equation template using $c_0=0$ yields $J(0)$, the desired original
instance. Both computation premises are admitted directly, and the
only transport source is the nonempty term $c_0$.

This is exactly the inherited three-node normalization in
\path{Research/certified_zero/tests/producer.py}, function
\texttt{expand\_computation}: one transport and two computation
premises. Its source is $c_0$, even when both defects are corrected
simultaneously. Passive zero occurrences in an update are left alone.
Thus all computation instances used in the mathematical proof can
be expanded before compiling to the source-restricted checker.

\subsection{Scope of the result}

The new statement is the admissibility of all omitted literal-zero
source transports. The native term/formula language, arithmetic
schemas, and target theory remain the inherited ones. This theorem
justifies replacing a two-way equality/Q major-pattern test by a
source-discriminated test in a suitable certificate protocol. It
does not assert a minimum proof-checker size, a shortest proof
expansion, or a new characterization of the primitive recursive
functions. The displayed proof is a written universal derivation;
executable certificates, when supplied, validate particular
instances rather than substitute for this argument.


\section{A deterministic calculus with auxiliary formulas}
\label{sec:composition}

The construction combines two restrictions that have different roles.  The
source of an equality elimination must not be the literal term $0$, and the
root of a paired matrix determines its introduction rule.  There is no
independent case/progress choice.  The interaction is resolved by allowing
auxiliary formulas in proofs.  Only the public formulas below are assigned
the ordinary conditional reading in the bi-interpretation theorem.  In
particular, we do not assert a round-trip theorem for every auxiliary
formula.

\subsection{Terms, public formulas, and proof formulas}

There is one numerical sort.  The terms are
\[
 t ::= x\mid 0\mid S(t)\mid I(n,a;u.h),
\]
where $u$ is bound only in $h$.  Substitution is capture-free.  Write
\[
 1:=S0,\qquad 2:=S1,\qquad C(v,a,b):=I(v,a;u.b),
\]
with $u$ fresh in the last abbreviation.  For an equation $E\equiv(s=t)$,
put
\[
 G_v E := \bigl(C(v,0,s)=C(v,0,t)\bigr).
\]

The public language $\G$ is the guarded language of the inherited package:
every equation is public; $\D_{x,y}B$ is public if $B$ is a compound public
formula; and an atomic matrix immediately under $\D_{x,y}$ must have the
literal shape
\begin{equation}
 B(x,y)\equiv I(y,0;u.h)=I(y,0;u.k).
 \tag{G}\label{eq:public-guard}
\end{equation}
The bodies $h,k$ may contain $x,y$, the state variable, and fixed
parameters.  This is a syntactic condition.  Every such matrix has a
provable zero instance.  The designated contradiction is $0=1$.

The auxiliary proof language $\U$ is larger:
\[
 A ::= (s=t)\mid\Q_y A\mid\D_{x,y}A.
\]
Its paired matrices need not satisfy \eqref{eq:public-guard}.  Contexts are
finite lists of reusable auxiliary assumptions.  Public theoremhood means
that a public conclusion has a derivation in this auxiliary calculus; it
does not require every intermediate formula to be public.  The same
definition applies to consequence from public assumptions.

\subsection{All primitive axioms and rules}

The only nonlogical axiom schemas, available under every context, are
\[
 I(0,a;u.h)=a,\qquad
 I(Sn,a;u.h)=h[I(n,a;u.h)/u].
 \tag{I0,IS}
\]
The hypothesis rule is $\Gamma\vdash A$ whenever $A\in\Gamma$.
The remaining rules are
\[
 \frac{\Gamma\vdash s=t\qquad\Gamma\vdash F(s)}
      {\Gamma\vdash F(t)}
 \quad(s\text{ is not literally }0),
 \tag{Tr$^{-}$}\label{eq:restricted-tr}
\]
\[
 \frac{\Gamma\vdash\Q_y B(y)\qquad\Gamma\vdash B(0)}
      {\Gamma\vdash B(t)},
 \tag{QE}
\]
\[
 \frac{\Gamma\vdash\D_{x,y}B(x,y)\qquad
       \Gamma,\Q_yB(a,y)\vdash A}
      {\Gamma\vdash A},
 \tag{DE}
\]
and one root-determined introduction schema
\begin{equation}
 \frac{\Gamma,B(t,\eta_B(z))\vdash B(t,Sz)}
      {\Gamma\vdash\D_{x,y}B(x,y)},\qquad
 \eta_B(z)=
 \begin{cases}
 z,&B\text{ is a literal equation},\\
 0,&B\text{ has root }\Q\text{ or }\D.
 \end{cases}
 \tag{DI$_{\rm root}$}\label{eq:root-di}
\end{equation}
Here $F$ is an auxiliary formula template.  In DE, $a$ is fresh for the
retained context, the conclusion, and the fixed parameters.  In DI,
$z$ is fresh for the retained context, the witness term $t$, and the
fixed parameters.  The displayed local assumption is discharged.
Bound names are always chosen to avoid capture.

There is no primitive equality reflexivity or successor axiom, no
separate induction schema, no unrestricted atomic case introduction,
and no primitive introduction rule for $\Q$.  Weakening, renaming,
contraction, and replacement of a retained assumption by its proof are
admissible by the usual induction on a finite derivation.  No unrestricted
numerical-substitution theorem for auxiliary derivations is needed below:
a substitution could turn an allowed source variable into the literal
$0$, so a direct constructor-wise substitution argument would be invalid.

\subsection{Induction and cases are available before equality is repaired}

\begin{lemma}[Auxiliary equation induction]\label{lem:aux-ei}
For every literal equation $E(y)$, the rule
\[
 \frac{\Gamma\vdash E(0)\qquad
       \Gamma,E(z)\vdash E(Sz)}
      {\Gamma\vdash E(t)}
\]
is admissible, with the usual fresh $z$.  The retained context may contain
arbitrary auxiliary formulas.
\end{lemma}
\begin{proof}
Apply atomic DI to $\D_{x,y}E(y)$, with $x$ unused and witness $0$.
Its premise is precisely the supplied step.  DE exports $\Q_yE(y)$;
the fresh witness is absent from this formula.  QE and the supplied base
give $E(t)$.  This derivation uses no transport at all.  Its paired
intermediate formula is generally auxiliary rather than public.
\end{proof}

\begin{lemma}[Numerical cases]\label{lem:aux-cases}
From $\Gamma\vdash A(0)$ and $\Gamma\vdash A(Sz)$, with $z$ fresh,
one may infer $\Gamma\vdash A(t)$ for every auxiliary $A(y)$.
\end{lemma}
\begin{proof}
For an equation, use Lemma~\ref{lem:aux-ei} and ignore its local induction
assumption.  For a compound formula, case DI with an unused witness
gives $\D_{x,y}A(y)$; DE and QE use the given base to obtain $A(t)$.
\end{proof}

This is the decisive bootstrap.  We do not first try to pad the hypothesis
$0=t$ into a public compound formula.  Auxiliary equation induction is
already available, and supplies the atomic cases needed by the
literal-zero-source argument.

These are the compatibility steps needed to combine the two previously
separate reductions.  The guarded public grammar is inherited from
\cite{guarded}; the source restriction is inherited from \cite{source}.

\subsection{Recovering equality from the restricted source rule}

The following proof reconstructs the zero-source argument in the pinned
\texttt{source\_dispatch} package and
\texttt{NativeZeroSource.lean}.  Its only non-equational ingredients are
equation induction and numerical cases, now supplied by
Lemmas~\ref{lem:aux-ei} and~\ref{lem:aux-cases}.  The formulas in the
calculations below are equations.
The new proof reuses this numerical argument only; the final unrestricted
auxiliary-formula transport argument of the source package is not used
\cite{source,lean-zero}.

\begin{lemma}[Safe equality calculations]\label{lem:safe-equality}
Reflexivity is derivable for every term.  Symmetry is derivable for $p=q$
when $p$ is not literally $0$, and transitivity is derivable through any
middle term that is not literally $0$.  Congruence and replacement on
either side of an equation are available along an equation with such a
source.  If $p,q$ are not literally $0$, then
\[
 \Gamma\vdash p=0,\quad\Gamma\vdash q=0
 \quad\Longrightarrow\quad\Gamma\vdash p=q.
 \tag{Zero join}\label{eq:zero-join}
\]
\end{lemma}
\begin{proof}
Put $r=I(0,t;u.0)$.  The axiom $r=t$, used as both premises of
Tr$^{-}$ with template $v=t$, gives $t=t$.  The usual equality
templates then give symmetry at an allowed source, transitivity through
an allowed middle term, and congruence at an allowed source.

For the last assertion, let $W(v)=I(0,v;u.0)$ and $w=W(0)$.  From
$p=0$ and $W(p)=W(p)$, replace $p$ only on the left to obtain
$w=W(p)$.  Symmetry at $w$, followed by the computation $W(p)=p$
on the left, gives $p=w$.  Similarly $q=w$.  Symmetry at $q$ and
transitivity through $w$ give $p=q$.  All these sources and the middle
term $w$ are syntactically nonzero; no numerical inequality is asserted.
\end{proof}

\begin{lemma}[Zero-source symmetry]\label{lem:zero-symmetry}
For every term $t$ and every auxiliary context,
\[
 \Gamma\vdash 0=t\quad\Longrightarrow\quad\Gamma\vdash t=0.
\]
\end{lemma}
\begin{proof}
The case where $t$ is literally $0$ is reflexivity.  Otherwise $t$ is an
allowed source.  Write
\[
 A(n)=I(n,t;u.Su).
\]
The zero computation and safe symmetry give $t=A(0)$, which may replace
the right side of the assumed $0=t$; thus $0=A(0)$.  Under
$z=A(z)$, congruence at the variable source $z$ gives
$Sz=S(A(z))$, and the successor computation gives $Sz=A(Sz)$.
Auxiliary equation induction therefore gives $1=A(1)$ and hence
\begin{equation}
 1=St.
 \label{eq:one-successor-t}
\end{equation}

For completeness, the small test used to remove this successor is
\[
 \begin{aligned}
 Z(v)&=C(v,1,0),& K(v)&=C(v,0,1),& H(v)&=C(v,1,2),\\
 N(n)&=I(n,0;u.Z(u)),& R(n)&=I(n,0;u.H(u)),& L(n)&=N(R(n)).
 \end{aligned}
\]
Numerical cases and computation give $H(v)=S(K(v))$.  Consequently
$R(1)=1$ and $R(SSv)=2$.  Computation gives $N(1)=1$ and $N(2)=0$,
so $L(1)=1$ and $L(SSv)=0$.  A second numerical case argument yields
\begin{equation}
 L(Sv)=Z(v).
 \label{eq:singleton-test}
\end{equation}
At the zero branch both sides reduce to $1$; at the successor branch
both sides reduce to $0$, and \eqref{eq:zero-join} joins them.  All
reductions use iterator sources.  Similarly numerical cases prove the
reversed filter identity
\begin{equation}
 C(Z(v),v,0)=v.
 \label{eq:reversed-filter}
\end{equation}
Its zero branch reduces via $Z(0)=1$, and its successor branch via
$Z(Sw)=0$.  Again every major used for reduction has a nonzero
syntactic source.

Applying $L$ to \eqref{eq:one-successor-t}, using
$L(1)=1$ and \eqref{eq:singleton-test}, and using safe symmetry and
transitivity, gives $Z(t)=1$.  Hence the left side
$q=C(Z(t),t,0)$ of \eqref{eq:reversed-filter} is also provably $0$.
The filter gives $q=t$.  Symmetry at the iterator source $q$, followed
by transitivity through $q$, yields $t=0$.

Every identity just used is a schema constructed at the displayed
terms.  The argument does not assume that arbitrary restricted
derivations may be numerically substituted.
\end{proof}

\begin{corollary}[Unrestricted equational transport]
Symmetry, transitivity, term congruence, and transport in every equation
template are admissible without a source restriction.
\end{corollary}
\begin{proof}
Symmetry follows from Lemmas~\ref{lem:safe-equality} and
\ref{lem:zero-symmetry}.  For transitivity through literal $0$, the
cases of a literal-zero endpoint are immediate; otherwise reverse the
second equation and use zero join.  For congruence along $0=t$ with
$t$ not literally $0$, reverse to $t=0$, use safe congruence, then
use unrestricted symmetry.  If $F(v)$ is $a(v)=b(v)$, term
congruence, symmetry, and transitivity transform $F(s)$ to $F(t)$.
\end{proof}

\begin{lemma}[Transport in public formulas]\label{lem:public-transport}
The unrestricted transport rule is admissible whenever its template $F$
is public.  Thus numerical equality has its full ordinary use throughout
the public language.
\end{lemma}
\begin{proof}
Induct on the public formula template.  The equation case is the preceding
corollary.  Suppose first that $F(v)=\D_{x,y}B(v,x,y)$ has compound
matrix.  Open $F(s)$ at a fresh witness $a$, obtaining the clause
$\Q_y B(s,a,y)$.  Introduce $F(t)$ at witness $a$ by compound DI.
Its base assumption $B(t,a,0)$ transfers, using the induction hypothesis
and $t=s$, to $B(s,a,0)$.  QE gives $B(s,a,Sz)$; the induction
hypothesis and $s=t$ transfer this to $B(t,a,Sz)$.

If $B$ is atomic, its public guard makes $B(s,a,0)$ a theorem.
The opened source clause therefore gives $B(s,a,Sz)$ outright.
Unrestricted equational transport gives $B(t,a,Sz)$.  Atomic DI
introduces the target, ignoring its local induction assumption.
Close DE in both cases.  Eigenvariables are chosen fresh for all fixed
parameters.
\end{proof}

This lemma is intentionally about public templates.  We do not need, and
do not assert, unrestricted zero-source transport in every auxiliary
formula containing an unguarded atomic paired matrix.

\subsection{Soundness of every auxiliary derivation}

Give auxiliary formulas the following ordinary first-order readings over
the scalar terms:
\[
 \begin{aligned}
 (s=t)^\dagger&=(s=t),\\
 (\Q_y B)^\dagger&=B^\dagger(0)\to\forall y\,B^\dagger(y),\\
 (\D_{x,y}B)^\dagger&=
 \begin{cases}
 \exists x\,\forall y\,(B(x,y)\to B(x,Sy)),&B\text{ is an equation},\\
 \exists x\,(B^\dagger(x,0)\to\forall y\,B^\dagger(x,y)),
       &B\text{ is compound}.
 \end{cases}
 \end{aligned}
\]
Numeric substitution preserves the root category of a formula, so this
reading commutes with capture-free numeric substitution.  The reading of
$\Q$ is always conditional, including at an equation.

\begin{theorem}[Auxiliary soundness]\label{thm:aux-soundness}
If $\Gamma\vdash_{\U}A$, then iPRA proves $A^\dagger$ from
$\Gamma^\dagger$, after the inherited primitive-recursive realization
of scalar terms.
\end{theorem}
\begin{proof}
Induct on the derivation.  Hypothesis, computation, transport, and QE
have their ordinary sound interpretations.  Atomic DI directly proves
the displayed progressive reading.  Compound DI uses numerical coverage
to extend its zero-to-successor premise to a zero-to-all conclusion.
For DE at an atomic matrix, ordinary equation induction converts
$\forall y(E(y)\to E(Sy))$ into $E(0)\to\forall yE(y)$, which is
exactly the exported $\Q$ assumption.  This uses only quantifier-free
induction, with the existential witness and other parameters fixed.
DE at a compound matrix is ordinary existential elimination.
All side assumptions and eigenvariable conditions are preserved.
\end{proof}

Let $A^\circ$ denote the \emph{uniform ordinary reading} on public
formulas: equality is actual numerical equality and every public
$\D_{x,y}B$ means $\exists x(B^\circ(x,0)\to\forall yB^\circ(x,y))$.

\begin{lemma}[Agreement on public formulas]\label{lem:public-agreement}
For every public $A$, iPRA proves $A^\dagger\leftrightarrow A^\circ$.
\end{lemma}
\begin{proof}
At a guarded atomic matrix $E(x,y)$, the zero instance $E(x,0)$ is a
theorem.  Progress implies all instances by equation induction, while
conditional closure implies all instances by that true base and hence
implies progress.  These two readings are therefore equivalent, with
$x$ fixed, and remain equivalent under the existential quantifier.
At compound matrices use structural induction.  Equations are unchanged.
\end{proof}

\subsection{Completeness and both interpretation composites}

We record the inherited guarded translation explicitly, because it
identifies the exact scope of the result.  Let $\mathcal C$ be the
existing native scalar calculus.  Define
\[
 L(E)=\D_{u,v}G_vE,\qquad
 P(A)=\begin{cases}L(A),&A\text{ an equation},\\A,&A\text{ compound}.
 \end{cases}
\]
The new variables in $L$ are fresh.  Translate old Eq/$\D$ formulas by
\[
 T(s=t)=(s=t),\qquad
 T(\D_{x,y}B)=\D_{x,y}P(TB).
\]
If the old presentation also treats $\Q$ as an object constructor, put
$T(\Q_yB)=\D_{x,y}P(TB)$ with $x$ unused.  Let $J$ be the literal
inclusion of public formulas into the old language.

\begin{lemma}[Guarded translation and its composites]\label{lem:guarded-composites}
The translation $T$ preserves derivations of the old calculus in the new
public consequence relation.  Uniformly under additional contexts,
\[
 A\dashv\vdash_{\mathcal C}JT(A),\qquad
 B\dashv\vdash_{\U}TJ(B)\quad(B\text{ public}).
\]
\end{lemma}
\begin{proof}
After equality repair, $L(E)\dashv\vdash E$: from $E$, atomic DI
proves $L(E)$ by ignoring its local assumption; conversely DE, the
theorem $G_0E$, and QE at $1$ recover $E$.  Therefore $P(A)$ is
always compound and interderivable with $A$.

Old case introduction transfers to compound DI after padding the base
and the successor conclusion.  Old equation progress transfers by
auxiliary equation induction under the padded base assumption, followed
by compound DI.  The old direct $\Q$ progress rule is covered by
Lemma~\ref{lem:aux-ei} and the unused-witness translation of $\Q$.
Elimination transfers by replacing each use of its opened clause with
padding of the antecedent, use of the translated clause, and unpadding
of the consequent.  Transport transfers by
Lemma~\ref{lem:public-transport}.  Computation and assumptions are
unchanged.  This is an effective transformation of finite proofs.

The old calculus has binder congruence and the same padding equivalence,
so structural induction gives its round trip.  In the new calculus,
uniform equivalence of two compound matrices lifts under $\D$ by DE
and compound DI.  At a guarded atomic matrix $B$, the remaining step is
\[
 \D_{x,y}B\dashv\vdash\D_{x,y}L(B).
\]
The forward direction opens the source clause, unpads the target base,
uses the clause, and pads the successor.  For the reverse direction,
the true guarded base proves $L(B(a,0))$; the opened clause then gives
$L(B(a,Sz))$ and hence $B(a,Sz)$.  Atomic DI introduces $\D B$,
ignoring its local assumption.  These are all the cases in the new
round-trip induction.  An unused existential witness introduces or
eliminates the old $\Q$ interpretation, using the witness $0$.
\end{proof}

\begin{theorem}[Exact public consequence relation]\label{thm:exact-public}
For public $\Gamma,A$,
\[
 \Gamma\vdash_{\U}A
 \quad\Longleftrightarrow\quad
 \Gamma\vdash_{\mathcal C}A
 \quad\Longleftrightarrow\quad
 \Gamma^\circ\vdash_{\ipra}A^\circ.
\]
\end{theorem}
\begin{proof}
Auxiliary soundness and public agreement give the rightward implication
from the new calculus to iPRA.  The pinned native--iPRA equivalence gives
the old native derivation.  Conversely, translate an old derivation by
$T$ and use both public round-trip equivalences, including the uniform
equivalences for the retained assumptions, to recover the original public
sequent.  The second equivalence is the inherited native--iPRA bridge.
\end{proof}

\begin{corollary}[Bi-interpretability and equiconsistency]
The ordinary realization of this public theory is bi-interpretable and
equiconsistent with full first-order iPRA.  Both numerical domains are the
whole domain; equality is actual equality; and both numerical composites
have the identity graph $x=y$.
\end{corollary}
\begin{proof}
Compose $T,J$ and their formula composites with the pinned
bi-interpretation of the native calculus and iPRA.  No term or numerical
sort is changed by $T$ or $J$, so the numerical composites remain the
identity.  The distinguished contradiction $0=1$ is unchanged and has
the ordinary reading equivalent to falsity in iPRA.  Exact public
consequence therefore gives equiconsistency as well.
\end{proof}

\subsection{Why the auxiliary distinction is essential}

\begin{proposition}[The canonical-only composition would fail]
If paired matrices throughout derivations were required to satisfy the
public grammar, with $\Q$ resources only the corresponding opened
clauses, the simultaneous restrictions would not even prove $1=0$ from
$0=1$.
\end{proposition}
\begin{proof}
Interpret terms in the standard natural numbers.  At the root of a
formula, declare every literal equation $0=t$ true, and interpret every
other equation by actual equality.  Interpret public compounds by their
ordinary conditional reading.  Every atomic leaf under a public paired
binder has an iterator on the left; consequently those leaves always
receive actual equality, and their readings commute with substitution.

Computations, guarded atomic progress, compound case introduction,
and the associated clause elimination are sound for this interpretation.
Restricted transport is sound too.  Its major $s=t$ has a syntactically
nonzero left side, so it asserts actual equality.  If an equation
template instantiated at $s$ has literal $0$ on the left, that left
template was constant $0$: it could not have been the hole, since
$s$ is not literally $0$.  The target is then also true by convention.
Otherwise transport preserves actual equality, unless its target has
literal-zero left side and is true by convention anyway.  Transport
inside a public compound is ordinary extensional transport.

The assumption $0=1$ is true in this interpretation, while $1=0$ is
false.  This proves the stated nonderivability.  The auxiliary induction
used in Lemma~\ref{lem:zero-symmetry} is exactly the extra proof resource
that invalidates this counterinterpretation.
\end{proof}

\begin{remark}[Exact checking contract]
The compact checker is a checker of auxiliary derivations with a
\emph{canonical public-goal contract}.  Its soundness and completeness
claim is Theorem~\ref{thm:exact-public}, for a goal in $\G$.  A raw
auxiliary goal need not be public, even when the program accepts its
proof.  The separately supplied strict checker includes and charges an
explicit public-goal validator.  A claim about the compact kernel alone
must not silently include the validator, and a claim about the whole
auxiliary syntax must not be substituted for the public bi-interpretation
proved here.
\end{remark}


\section{A uniform arity encoding and its checking theorem}
\label{sec:arity-checker}

The arithmetic and object logic remain the scalar-iteration Eq/$\D$
calculus specified above.  The change is in its finite certificate
representation.  A paired formula is unary, an equation is binary, and
an internal closure is ternary:
\[
 \D_{x,y}B\longleftrightarrow[B],\qquad
 s=t\longleftrightarrow[s,t],\qquad
 \Q_yB\longleftrightarrow[B,B,B].
 \tag{Arity encoding}\label{eq:arity-encoding}
\]
The three copies in a closure are representation redundancy: only the
third has mathematical content.  Their benefit is that tripling a unary
major constructs exactly its closure.  Tripling any other nonempty
formula array produces an array of arity at least six, which will denote
truth.  This makes the other major cases neutral without a separate
test.

The previous generated-assumption and constructed-computation interfaces
are inherited from the corpus~\cite{generated,constructed}.  The new
encoding makes a single arithmetic expression, $5-\operatorname{arity}$,
perform the necessary depth adjustments.  Its soundness depends on an
explicit interpretation of surplus syntax.  It would not be sound to
drop arity checks while continuing to treat every surplus array as an
unconstrained new connective.

\subsection{The complete sorted interpretation of arrays}
\label{subsec:sorted-arrays}

Write $\top$ for the fixed equation $0=0$.  A protected numerical slot
has width two.  The numerical root has depth three: the coordinates
$0,1,2$ all denote zero, and under $p$ protected slots the depth is
$3+2p$.  Coordinates $3+2j$ and $4+2j$ name slot $j$, counted from
the innermost slot, for $0\leq j<p$.  Formula depth is $2p$ under the
same $p$ slots.  Negative coordinates admitted by the current global
floor name rigid numerical parameters.

The decoder is sorted by the position in which a tree is used.  Let
$\rho$ be its ordered environment of protected numerical slots.
The notation $(a,\rho)$ means prepend a slot of value $a$.
A superscript $\sharp$ denotes a child decoded in the indicated
environment.

\begin{center}
\begin{tabular}{cll}
\toprule
Numerical arity & Meaning & Relevant child environment\\
\midrule
$0$ & $0$ & No child\\
$1$: $[b]$ & $Sb^\sharp$ & $(0,0,\rho)$\\
$2$ & $0$ & Annotations ignored\\
$3$: $[n,a,h]$ & $I(n^\sharp,a^\sharp;v.h^\sharp)$ & $(0,\rho)$ for $n,a$; $(v,\rho)$ for $h$\\
$\geq4$ & $0$ & All children ignored\\
\bottomrule
\end{tabular}
\end{center}

\begin{center}
\begin{tabular}{cll}
\toprule
Formula arity & Meaning & Relevant child environment\\
\midrule
$0$ & $\top$ & No child\\
$1$: $[B]$ & $\D_{x,y}B^\sharp$ & $(y,x,\rho)$\\
$2$: $[s,t]$ & $s^\sharp=t^\sharp$ & $\rho$ numerically\\
$3$: $[P,U,B]$ & $\Q_yB^\sharp$ & $(y,\rho)$ for $B$\\
$\geq4$ & $\top$ & All children ignored\\
\bottomrule
\end{tabular}
\end{center}

The first two entries of a ternary formula are annotations.  A canonical
closure copies its body into both of them, as in
\eqref{eq:arity-encoding}.  Successfully admitted scalar formula padding
also denotes $\top$.  An internal closure has the ordinary meaning
$B(0)\to\forall yB(y)$; it is not an equation and is not eligible for
compound induction.

Canonical encoded object expressions use arrays and integer coordinates,
respect these slots, and contain no high-arity node.  The checker itself
still generates high-arity truth entries in administrative contexts,
even during a canonical proof.  In particular, constructing a
successor beneath positive bound coordinates requires passing through
two dummy zero slots.  Constructing an iterator requires one new slot:
zero in its counter and seed, and its genuine state variable in the
update.  Rigid negative parameter names are unaffected by these shifts.

The additional raw shapes are certificate aliases with the displayed
meanings.  They are not independent predicate or function symbols added
to the object theory.  Our exact completeness statement concerns the
canonical native language.  We do not require every possible raw spelling
of a semantic alias to be interchangeable under the same raw certificate.

\subsection{One depth calculation replaces the lookup table}

For every array the opener recursively adds $5-\operatorname{arity}$
to its depth.  The meaningful cases are
\[
\begin{array}{c|ccc}
 \operatorname{arity}&1&2&3\\\hline
 \text{depth increment}&4&3&2.
\end{array}
\]
Unary $\D$ adds two numerical binders; unary successor adds two dummy
zero slots.  Binary equality changes formula depth $2p$ to numerical
depth $3+2p$.  Ternary $\Q$ adds one numerical binder; ternary iteration
adds one state slot, or one dummy slot in its counter and initializer.
The same calculation therefore covers all relevant constructors.

Binary numerical annotations denote zero, so their parity-changing
traversal has no mathematical effect.  Empty arrays have no children.
At arity at least four, the entire node already denotes a constant.
Even a zero or negative depth increment below such a node cannot expose
one of its ignored children as a meaningful variable.  Traversal may
still reject an alias; it cannot change its stipulated interpretation.

This is a semantic assertion, not raw idempotence of the opener.
For example, at floor $-1$, opening a six-element array of zeros with
image $-1$ can yield six copies of $-1$, while a later no-image opening
of that result fails at a negative child depth.  Both formula arrays
denote $\top$.  The proof below only concerns successful calls at meaningful
sorted positions.  Canonical input object expressions contain no surplus
arrays; the administrative exports that generate them are justified
separately as truth entries.

\subsection{Opening and scope}
\label{subsec:opening-lemma}

Write $i(a,t,d)$ for the executable opener at floor $k$.  It preserves a
scalar when $k\leq a<d$.  At a numerical position, every other relevant
scalar is the distinguished substitution hole; the supplied image is
checked at numerical root depth three before insertion.  At a formula
position scalar padding is truth and carries no numerical hole.
Missing images at used numerical holes reject.

\begin{lemma}[Semantic opening]
\label{lem:semantic-opening}
For every successful call at a meaningful formula root of depth $2p$
or numerical root of depth $3+2p$, with $p\geq0$, decoding its result
is capture-free filling of the distinguished hole in the decoded input
template.  These depths protect $p$ numerical slots; admitted rigid
parameters are preserved.  If an image is semantically unused, zero
may be chosen as its mathematical representative.  Recursive calls
at negative depths inside ignored annotations need no independent
sorted interpretation.
\end{lemma}

\begin{proof}
Induct on the finite raw tree.  At a relevant scalar, the interval
comparison is exactly the protected-coordinate test.  At a meaningful
unary, binary, or ternary node, the preceding tables supply precisely
the required additional slots.  At numerical binary annotations,
high-arity nodes, and ignored closure fields, the relevant mathematical
meaning is constant or absent, so the conclusion is immediate
regardless of traversal inside the annotation.

An inserted image is checked at its own numerical root.  Its free
parameters are rigid negative names, and it has no unbound positive
coordinate there.  Insertion beneath protected slots therefore captures
none of its parameters; its internal binders retain their relative
indices.  This proves the inductive step for substitution as well as
for mere traversal.
\end{proof}

A successful no-image check consequently preserves the interpretation
of its checked closed expression.  Raw structural equality of checked
outputs implies equality of their decoded meanings.  The checker does
not infer equality merely from a semantic comparison of two different
programs or numerical terms.

The global floor decreases by one at each non-hypothesis continuation.
If its old value is $K$, the integer $K-1$ is the designated fresh
eigenparameter.  With integer coordinates it is the only new name.
If fractional coordinates are accepted, other names in $[K-1,K)$ may
also occur in the finite certificate.  These are additional fresh
parameters.  They must not be silently equated to the designated
integer eigenparameter.

Previously checked relevant goals and retained hypotheses mention none
of these new names: their earlier no-image checks would reject a
relevant numerical occurrence below the old floor.  Coordinates inside
ignored annotations do not contribute free variables to the decoded
formula.  This is the freshness invariant used in every rule below.

\subsection{A computation hypothesis only needs a zero instance}
\label{subsec:zero-supply}

The shared continuation can open a generated arithmetic equation at a
fresh parameter.  That equation need not be a theorem uniformly in the
new parameter.  The following admissible rule is sufficient.

\begin{lemma}[Supply with a known zero instance]
\label{lem:zero-supply}
Suppose $\vec a$ is fresh for $\Gamma,F$, the calculus proves
$H(\vec0)$, and $\Gamma,H(\vec a)\vdash F$.  Then
$\Gamma\vdash F$.
\end{lemma}

\begin{proof}
Substitute zero for every member of $\vec a$ in the second derivation.
Freshness leaves $\Gamma,F$ unchanged.  Cut in the derivation of
$H(\vec0)$.  Both substitution and cut are admissible in the native
calculus.  This is a finite proof transformation and does not use a
consistency assumption.
\end{proof}

All names newly admitted when the floor changes belong to $\vec a$;
for noninteger aliases this includes the finite collection in
$[K-1,K)$ that actually occurs.  On decoding, their simultaneous zero
substitution agrees with opening the generated expression at zero using
the old floor.  Previously admitted parameters are preserved.  New
rigid names and unprotected placeholders both become zero.  The opener
may keep different raw zero spellings, but their decoded values agree.

The constructed-computation packet supplies a triple $s=[n,b,r]$ and a
mode $u$.  In successor mode its generated equation has raw form
\[
 \bigl[\,[[n],b,r],\ i(r,s,3)\,\bigr],
 \tag{Step supply}\label{eq:step-supply}
\]
while in zero mode it has form
\[
                    \bigl[\,[u,b,r],\ b\,\bigr].
 \tag{Zero supply}\label{eq:base-supply}
\]
Here successor mode means that $u$ is truthy; the zero mode has falsy
$u$.  These expressions describe actual arrays, before their final
opening.  In \eqref{eq:step-supply} the counter is constructed as a
successor, rather than accepted on an unsupported assertion that it is
one.  This construction is inherited from the corresponding draft
interface~\cite{constructed}.

\begin{lemma}[Sound constructed computation]
\label{lem:constructed-supply}
Every successful arithmetic supply has a decoded zero instance
derivable from the two scalar-iteration computation schemas.
\end{lemma}

\begin{proof}
First apply the simultaneous zero substitution described above.  In
zero mode every successful falsy numerical counter denotes zero, so
\eqref{eq:base-supply} is an $I0$ instance.  In the exact-list profile
the relevant accepted aliases are an empty array, numeric zero, false,
and numerically equal zero values.  Null, strings, and dictionaries
cannot enter a meaningful numerical position through an unchecked
coercion.

In successor mode the left counter is an actual successor.  If the
update uses its state in a semantically relevant position, opening it
invokes the no-image check of the recursive source $s$ at numerical
root three.  This checks the entire source, including its update.
Only coordinates three and four are aliases for its single state slot.
A genuinely escaping coordinate in that update is revisited during
the recursive-image check and rejects.  Its counter and initializer
are likewise closed with their prescribed dummy zero slot.

Inserting the extra dummy slots of the displayed successor does not
alter such a checked predecessor: extra unbound slots have value zero,
and actual free parameters are rigid.  Lemma~\ref{lem:semantic-opening}
therefore identifies the right side of \eqref{eq:step-supply} with the
update applied to exactly the recursive predecessor value occurring
on its left.  This is the $IS$ schema.

If the update has no semantically relevant state occurrence, its
decoded value is constant in that variable.  Its zero-normalized value
is the same in the source update and on the right, independently of
the unused recursive image; this is again an $IS$ instance.  Traversal
inside an ignored annotation may still inspect the image and reject.
That removes a certificate and does not invalidate the argument for a
successful one.

Finally, successful three-field source extraction cannot hide another
JSON operator.  A scalar is not iterable.  A three-character string,
or a three-key JSON dictionary, necessarily puts a nonempty string in
a relevant source field, which the inspected opener rejects.  The
successful source is therefore the stated array triple.
\end{proof}

The recursive-image check in the state-dependent case is essential.
Merely opening a final displayed equation at zero would not establish
that its two state interpretations agree.  Combining
Lemma~\ref{lem:constructed-supply} with Lemma~\ref{lem:zero-supply}
justifies both actual image choices of the common continuation, including
the branch that retains a fresh-parameter hypothesis.

\subsection{Audit of the logical continuations}
\label{subsec:logical-audit}

\paragraph{Hypothesis closure.}
The initial opener check and membership comparison select the same
decoded context formula.  Additional aliases can make a spelling less
convenient, but cannot identify an equation with a false denotation
solely by semantic evaluation.

\paragraph{Equality transport.}
The major must match $[s,u]$, and opening the proposed formula template
at $u$ must match the target.  The continuation checks its opening at
$s$.  Lemma~\ref{lem:semantic-opening} gives ordinary transport along
the decoded equality.  The argument applies to compound formulas,
internal clauses, and truth aliases as well as equations.

\paragraph{Closure specialization.}
The closure pattern is the triplicate $[t,t,t]$.  The guard forces the
old instance to have numerical image zero; the new instance may have
an arbitrary numerical image.  This is exactly
$B(0)\to\forall yB(y)$.  The zero guard is indispensable.  Without
it, a closure with false zero instance and true successor instances
would turn a proved successor instance into falsity.

\paragraph{Paired elimination.}
The export is always three concatenated copies of the major.  A unary
major $[B]$ therefore exports $[B,B,B]$.  All three copies are opened
identically; the third copy has precisely the witness instance required
by paired elimination.  The newly allocated witness is fresh for the
retained context and conclusion.

An equation major instead exports arity six; a closure major exports
arity nine; a larger array exports still larger arity.  These are all
truth aliases.  Thus failed elimination endpoint claims have only the
effect of adjoining a theorem.  A numeric scalar major, if accepted as
formula padding, exports scalar padding after multiplication and also
denotes truth.  Unsupported types reject.

A ternary formula with unequal annotations still denotes the closure
of its third child.  It fails the triplicate specialization pattern
unless the three fields match.  This loses a possible raw spelling,
not a mathematical soundness condition.  Canonical compilation uses
triplicate closures throughout.

\paragraph{Paired introduction.}
The call \texttt{i(*a,t,2)} uses the function's argument count to enforce
the relevant outer arity.  For a unary goal $a=[B]$, it is exactly
\texttt{i(B,t,2)}: the universal slot is protected while the
existential witness is opened.  For arity at least two, there are too
many arguments and the call rejects.  An empty goal can instead make
the legal two-argument call \texttt{i(t,2)}.  Its original target was
already $\top$, so accepting it cannot add a false theorem.

For the unary goal, the subsequent successor opening gives the intended
premise $B(t,Sz)$.  The local hypothesis is $B(t,0)$ in case mode and
$B(t,z)$ in progress mode.  The freshness invariant supplies the
usual eigenvariable conditions.

\paragraph{Induction remains restricted.}
Among meaningful nonconstant matrix constructors, only binary equality
has even arity.  Unary paired formulas and ternary closures have odd
arity and receive case introduction.  Empty or surplus arrays that
reach the progress branch denote the fixed equation $0=0$, whose
induction is valid.  Scalar padding that cannot provide a length
rejects.  The broader certificate alphabet therefore adds no induction
on a quantified matrix.

\paragraph{Additional fresh parameters.}
An accepted subderivation can contain newly admitted fractional names,
but the older context and target do not depend on them.  Replace these
auxiliary names by zero when closing the rule, retaining its designated
integer eigenvariable where that rule needs one.  The common opening
of an already checked retained goal leaves its decoded meaning
unchanged.  This handles the actual global-floor convention without
pretending that every admissible negative number is an integer.

\subsection{Canonical completeness, including zero witnesses}
\label{subsec:canonical-completeness}

Rename a finite native proof's eigenvariables to fresh negative integers.
Encode its terms with the slot conventions of
Section~\ref{subsec:sorted-arrays}; encode paired formulas as unary and
closures as triplicate.  Hypothesis leaves consume no instruction.
Transport, closure specialization, and paired elimination have their
displayed packets.  Arithmetic computation leaves use
\eqref{eq:base-supply} or \eqref{eq:step-supply} and then close the
target by hypothesis lookup.

One representational issue requires an actual proof transformation.
The packet's truthy witness field selects introduction.  A literal zero
witness therefore cannot simply be inserted into that field.  Define
\[
                 J(t)=I(0,t;v.0).
\]
The computation axiom proves $J(t)=t$, and $J(t)$ has a nonempty array
representation even when its value is zero.  To introduce with a falsy
witness spelling $t$, introduce instead with $J(t)$: transport its
retained assumption to the assumption using $t$, use the original
premise derivation, and transport the resulting successor instance
back to the instance using $J(t)$.  Derived symmetry supplies the
reverse transport.  This is a finite transformation of the premise,
and the paired conclusion itself is unchanged.  If the existential
coordinate is unused, any closed truthy witness works without a
conversion.

The scalar coordinate $1$ is also a semantic zero at numerical root
three, but that observation alone is not a completeness proof.  Closure
specialization and case introduction insert literal raw zero, and
distinct zero aliases need not be interchangeable under a fixed raw
proof spelling.  The computation-backed box gives the required
derivation instead of assuming representational interchangeability.

The mode zero selects progress at an equation.  A different falsy
mode, such as the empty array, selects case.  These are certificate
control tags, not numerical assertions about the witness.  All steps
in the compilation are finite and effective.

\begin{theorem}[Conservativity and completeness of the arity checker]
\label{thm:arity-checker}
Every accepted certificate has its decoded conclusion provable in the
ordinary realization of the native calculus.  For every canonical
native end sequent, the internal checking judgment and the native
calculus have the same consequence relation, in the idealized
unbounded integer-and-array implementation.  In particular, the public
empty-context driver recognizes exactly the canonical closed theorems.
The translation in each direction is effective on finite derivations.
\end{theorem}

\begin{proof}
For soundness, induct on the finite accepting call tree.  Hypothesis
closure uses Lemma~\ref{lem:semantic-opening}; the logical cases are
the audited continuations above.  Arithmetic supply uses
Lemmas~\ref{lem:constructed-supply} and \ref{lem:zero-supply}.
Every discarded annotation is accounted for by its zero or truth
interpretation, and every fresh parameter by the scope invariant.
For completeness, compile a finite native derivation by the preceding
construction, including the boxed-witness transformation.  Its
canonical input expressions require no surplus-arity aliases or fractional
names; generated administrative surplus arrays have exactly the truth
meaning already proved.  This proves both implications without using a finite test
suite as a substitute for induction on derivations.
\end{proof}

The unchanged native calculus is bi-interpretable with full first-order
$\ipra$ by the inherited arithmetic and logical bridge
\cite{generated,corpus}.  In particular,
\[
 (\D_{x,y}B)^\circ
 =\exists x\bigl(B^\circ(x,0)\to\forall yB^\circ(x,y)\bigr).
\]
The whole numerical domain and actual equality are retained.  Its
internally proved numerical round trips and its formula-composite
derivations therefore transfer through
Theorem~\ref{thm:arity-checker}.  We have changed the checked
representation, not replaced bi-interpretability by standard-model
agreement or supplied a new proof of the scalar predecessor theorem.

\subsection{JSON and finite-runtime qualifications}

The primary profile uses an ordinary list-type test.  The smaller byte
profile instead tests multiplication by zero against an empty list.
Among standard JSON runtime values the latter recognizes exactly lists:
numeric values and booleans produce numeric zero, strings produce an
empty string, and null and dictionaries raise rather than becoming
arrays.  This argument assumes ordinary decoded JSON values, not
arbitrary Python objects with overloaded methods.

Finite nonnegative real coordinates have the same interval
interpretation: $[0,3)$ is zero, and $[3+2j,5+2j)$ names slot $j$.
Finite negative coordinates in the admitted range are additional rigid
names.  Equal Python numerical values have the same decoded role,
including Boolean aliases and signed zero.  Nonfinite values admitted
by Python's permissive JSON reader cannot survive a relevant closed
numerical check.  They may be filled as template holes or ignored
inside a constant annotation; neither makes them unchecked parameters.

The alternative opener based on generic iteration has a separate
runtime qualification.  Empty strings and empty dictionaries traverse
as empty arrays; a dictionary whose sole key is empty traverses as a
unary array with an empty child.  Dictionary values are not visited.
Nonempty strings can rely on CPython's recursion limit to reject, as
also occurs in a current corpus comparison profile~\cite{transient}.
The idealized structural-termination claim for the exact-list profile
must not be transferred to that generic profile unchanged.

All executable profiles have ordinary finite memory, stack, and input
limits.  These can reject a sufficiently large otherwise valid proof.
The formal consequence relation is the unbounded finite-certificate
calculus; the delivered program checks certificates within its runtime
resources.  The measurements charge the entire delivered checker,
while the semantic proof and canonical producer are supplied separately.


\clearpage
\section{Complete new checker listings}
\label{sec:listings}

Each listing is a complete acceptance module. The executable files in
\texttt{code/} are authoritative for the measurements and can be run
independently. The expanded explanatory implementation is additional
documentation, outside every measured trusted module.

\Needspace{8\baselineskip}
\subsection{Inferred-major token profile}
\begin{lstlisting}[language=Python]
from json import load
k = 0

def i(a, t=0j, d=0):
    try:
        return [i(b, t, d + 5 - len(a)) for b in a]
    except:
        return a if k <= a < d else d % 2 and i(t, d=3)

def c(a, g):
    global k
    while i(a) not in g:
        m, s, u, t = q()
        if i(t, u) == a:
            m = [s, u] if s else [t] * 3
            a = i(t, s)
        h = c(m, g) or m * 3 if m else [[s and [s], u, t], i(t, [s, u, t], 3) if s else u] if u else (a := i(*a, t, 2))
        k -= 1
        g += (i(h, (m or len(a) == 2) and k),)
        a = i(a, [k])
q = load(open(0)).pop
c(q(), ())
print('ACCEPT')
\end{lstlisting}
\Needspace{8\baselineskip}
\subsection{Inferred-major AST profile}
\begin{lstlisting}[language=Python]
from json import load

def c(a, g=(0,), k=0):

    def i(a, t=None, d=0):
        try:
            return [i(b, t, d + 5 - len(a)) for b in a]
        except:
            return a if k <= a < d else d % 2 and i(t, d=3)
    while i(a) not in g:
        m, s, u, t = q()
        if i(t, u) == a:
            m = [s, u] if s else [t] * 3
            a = i(t, s)
        h = c(m, g, k) or m * 3 or ([[s and [s], u, t], i(t, [s, u, t], 3) if s else u] if u else (a := i(*a, t, 2)))
        k -= 1
        g += (i(h, (m or len(a) == 2) and k),)
        a = i(a, [k])
c((q := load(open(0)).pop)())
print('ACCEPT')
\end{lstlisting}
\Needspace{8\baselineskip}
\subsection{Truthy-result byte profile}
\begin{lstlisting}[language=Python]
import json

def c(a, g=[0], k=0):
    i = lambda a, t=c, d=0: [i(b, t, d + 5 - len(a)) for b in a] if a * 0 != 0 else a if k <= a < d else d % 2 and i(t, d=3)
    return i(a) in g or [((h := (c(m, g, k) and m * 3 or (u and [[s and [s], u, t], i(t, [s, u, t], 3) if s else u]) or (a := i(*a, t, 2)))), c(i((m, a) == (s and [s, u] or [t] * 3, i(t, u)) and i(t, s) or a, [(k := (~-k))]), g + [i(h, (m or ~len(a) % 2) and k)], k)) for m, s, u, t in [q()]]
q = json.load(open(0)).pop
c(q())
print('ACCEPT')
\end{lstlisting}
\Needspace{8\baselineskip}
\subsection{Byte profile, ordinary decrement}
\begin{lstlisting}[language=Python]
import json

def c(a, g=[0], k=0):
    i = lambda a, t=c, d=0: [i(b, t, d + 5 - len(a)) for b in a] if a * 0 != 0 else a if k <= a < d else d % 2 and i(t, d=3)
    return i(a) in g or [((h := (c(m, g, k) and m * 3 or (u and [[s and [s], u, t], i(t, [s, u, t], 3) if s else u]) or (a := i(*a, t, 2)))), c(i((m, a) == (s and [s, u] or [t] * 3, i(t, u)) and i(t, s) or a, [(k := (k - 1))]), g + [i(h, (m or ~len(a) % 2) and k)], k)) for m, s, u, t in [q()]]
q = json.load(open(0)).pop
c(q())
print('ACCEPT')
\end{lstlisting}
\Needspace{8\baselineskip}
\subsection{Native public-goal profile}
\begin{lstlisting}[language=Python]
import json

def c(a, g=[0], k=0, v=0):
    i = lambda a, t=c, d=0: [i(b, t, d + 5 - len(a)) for b in a] if a * 0 == [] and v <= (len(a) in (1, 2 + d % 2)) else a if k <= a < d and v <= (type(a) is int and d % 2 and (not a or a > 1 == a % 2)) else d % 2 | v and i(t, d=3)
    if i(a) in g:
        return
    v = 0
    m, s, u, t = q()
    h = c(m, g, k) or m * 3 or (u and [[s and [s], u, t], i(t, [s, u, t], 3) if s else u]) or (a := i(*a, t, 2))
    c(i(m == ([s, u] if s else [t] * 3) and i(t, u) == a and i(t, s) or a, [(k := (k - 1))]), g + [i(h, (m or len(a) % 2 == u) and k)], k)
q = json.load(open(0)).pop
c(q(), v=1)
print('ACCEPT')
\end{lstlisting}

\begin{thebibliography}{99}
\bibitem{upstream}ProveIt contributors. Concise presentations of intuitionistic primitive recursive arithmetic. Pinned snapshot, 12 November 2026.
\href{https://github.com/VladimirReshetnikov/ProveIt/tree/62396b06020f794fa6dc01b4025a6a5db9d4d28a/Logic/PrimitiveRecursiveArithmetic/ConciseIPRA}{Pinned primary source}.
\bibitem{corpus}ProveIt contributors. Collected research packages and current formalization.
\href{https://github.com/VladimirReshetnikov/ProveIt/tree/62396b06020f794fa6dc01b4025a6a5db9d4d28a/Logic/PrimitiveRecursiveArithmetic/ConciseIPRA/Research}{Pinned primary source}.
\bibitem{generated}ProveIt contributors. Inherited arity encoding and generated-assumption proof; reproduced in this report.
\href{https://github.com/VladimirReshetnikov/ProveIt/tree/62396b06020f794fa6dc01b4025a6a5db9d4d28a/Logic/PrimitiveRecursiveArithmetic/ConciseIPRA/Research/checked_refinement/reference/foundations}{Pinned primary source}.
\bibitem{constructed}ProveIt contributors. Constructed-computation certificate interface.
\href{https://github.com/VladimirReshetnikov/ProveIt/tree/62396b06020f794fa6dc01b4025a6a5db9d4d28a/Logic/PrimitiveRecursiveArithmetic/ConciseIPRA/Research/constructed_computation}{Pinned primary source}.
\bibitem{guarded}ProveIt contributors. Guarded core and deterministic introduction.
\href{https://github.com/VladimirReshetnikov/ProveIt/tree/62396b06020f794fa6dc01b4025a6a5db9d4d28a/Logic/PrimitiveRecursiveArithmetic/ConciseIPRA/Research/guarded_core}{Pinned primary source}.
\bibitem{source}ProveIt contributors. Literal-zero-source admissibility and source dispatch.
\href{https://github.com/VladimirReshetnikov/ProveIt/tree/62396b06020f794fa6dc01b4025a6a5db9d4d28a/Logic/PrimitiveRecursiveArithmetic/ConciseIPRA/Research/source_dispatch}{Pinned primary source}.
\bibitem{lean-zero}ProveIt contributors. Lean formalization and source-admissibility results, as recorded in the pinned README.
\href{https://github.com/VladimirReshetnikov/ProveIt/tree/62396b06020f794fa6dc01b4025a6a5db9d4d28a/Logic/PrimitiveRecursiveArithmetic/ConciseIPRA}{Pinned primary source}.
\bibitem{erasure}ProveIt contributors. Inherited scalar predecessor bootstrap; exact attributed appendix reproduced in this report.
\href{https://github.com/VladimirReshetnikov/ProveIt/tree/62396b06020f794fa6dc01b4025a6a5db9d4d28a/Logic/PrimitiveRecursiveArithmetic/ConciseIPRA/Research/checked_refinement/reference/foundations}{Pinned primary source}.
\bibitem{progress}ProveIt contributors. Inherited equality reflection, pairing, and primitive-recursion compilation appendices; reproduced in this report.
\href{https://github.com/VladimirReshetnikov/ProveIt/tree/62396b06020f794fa6dc01b4025a6a5db9d4d28a/Logic/PrimitiveRecursiveArithmetic/ConciseIPRA/Research/checked_refinement/reference/foundations}{Pinned primary source}.
\bibitem{forward}ProveIt contributors. Inherited native arithmetic and proof-certificate foundations.
\href{https://github.com/VladimirReshetnikov/ProveIt/tree/62396b06020f794fa6dc01b4025a6a5db9d4d28a/Logic/PrimitiveRecursiveArithmetic/ConciseIPRA/Research/checked_refinement/reference/foundations}{Pinned primary source}.
\bibitem{transient}ProveIt contributors. Transient-seed arithmetic development.
\href{https://github.com/VladimirReshetnikov/ProveIt/tree/62396b06020f794fa6dc01b4025a6a5db9d4d28a/Logic/PrimitiveRecursiveArithmetic/ConciseIPRA/Research/transient_seed}{Pinned primary source}.
\end{thebibliography}

\end{document}
